\documentclass[11pt,a4paper]{article}

\usepackage[a4paper,margin=20mm]{geometry}
\usepackage{amsmath,amssymb}
\usepackage{booktabs,array}
\usepackage[hidelinks]{hyperref}

\setlength{\parskip}{6pt plus 2pt minus 1pt}
\setlength{\parindent}{0pt}

\newcommand{\optsfourA}[4]{%
  \begin{tabbing}
    \hspace{0.48\linewidth}\=\kill
    (A) #1 \> (B) #2 \\[2pt]
    (C) #3 \> (D) #4
  \end{tabbing}%
}

\newcommand{\optslistA}[4]{%
  \begin{itemize}
    \setlength{\itemsep}{2pt}\setlength{\parskip}{0pt}
    \item[(A)] #1
    \item[(B)] #2
    \item[(C)] #3
    \item[(D)] #4
  \end{itemize}%
}

\newcommand{\ans}[1]{\textbf{Answer:} #1\par\medskip}

\begin{document}

% ============================================================================
% GATE 2012 -- MATHEMATICS (MA) & GENERAL APTITUDE (GA)
% ============================================================================

\begin{center}
  {\Large\bfseries GATE 2012 -- Mathematics (MA)}\\[4pt]
  {\small Duration: Three Hours \quad\textbar\quad Maximum Marks: 100}
\end{center}

\medskip
\textbf{Notations:} $\mathbb{R}, \mathbb{C}, \mathbb{Z}$: real, complex, integer sets; $F$: field; $\mathbb{C}^n, F^n$: $n$-tuples; $N(\mu, \sigma^2)$: normal distribution; $E(X), \operatorname{Cov}(X,Y)$: expectation and covariance; $S_n$: permutation group on $n$ symbols; $P_n$: polynomials of degree at most $n$; $C_n$: cyclic group of order $n$; $Z(G)$: centre of $G$; $i = \sqrt{-1}$.

\section*{Q.1 -- Q.25 Carry One Mark Each}

\textbf{Q.1} The straight lines $L_1 : x = 0$, $L_2 : y = 0$, and $L_3 : x + y = 1$ are mapped by the transformation $w = z^2$ into the curves $C_1, C_2$, and $C_3$ respectively. The angle of intersection between the curves at $w = 0$ is
\optsfourA{$0$}{$\pi/4$}{$\pi/2$}{$\pi$}
\ans{(D)}

\textbf{Q.2} In a topological space, which of the following statements is NOT always true:
\optslistA{Union of any finite family of compact sets is compact.}{Union of any family of closed sets is closed.}{Union of any family of connected sets having a non-empty intersection is connected.}{Union of any family of dense subsets is dense.}
\ans{(B)}

\textbf{Q.3} Consider the following statements:\\
$P$: The family of subsets $\left\{A_n = \left(-\dfrac{1}{n}, \dfrac{1}{n}\right),\ n = 1, 2, \dots\right\}$ satisfies the finite intersection property.\\
$Q$: On an infinite set $X$, a metric $d : X \times X \to \mathbb{R}$ is defined as $d(x, y) = \begin{cases} 0, & x = y, \\ 1, & x \ne y. \end{cases}$ The metric space $(X, d)$ is compact.\\
$R$: In a Fr\'echet ($T_1$) topological space, every finite set is closed.\\
$S$: If $f : \mathbb{R} \to X$ is continuous, where $\mathbb{R}$ is given the usual topology and $(X, \tau)$ is a Hausdorff ($T_2$) space, then $f$ is a one-one function.\\
Which of the above statements are correct?
\optsfourA{$P$ and $R$}{$P$ and $S$}{$R$ and $S$}{$Q$ and $S$}
\ans{(A)}

\textbf{Q.4} Let $H$ be a Hilbert space and $S^\perp$ denote the orthogonal complement of a set $S \subseteq H$. Which of the following is INCORRECT?
\optsfourA{For $S_1, S_2 \subseteq H$; $S_1 \subseteq S_2 \Rightarrow S_1^\perp \subseteq S_2^\perp$}{$S \subseteq (S^\perp)^\perp$}{$\{0\}^\perp = H$}{$S^\perp$ is always closed}
\ans{(A)}

\textbf{Q.5} Let $H$ be a complex Hilbert space, $T : H \to H$ be a bounded linear operator, and let $T^*$ denote the adjoint of $T$. Which of the following statements are always TRUE?\\
$P$: $\forall x, y \in H,\ \langle Tx, y \rangle = \langle x, T^* y \rangle$\\
$Q$: $\forall x, y \in H,\ \langle x, Ty \rangle = \langle T^* x, y \rangle$\\
$R$: $\forall x, y \in H,\ \langle x, Ty \rangle = \langle x, T^* y \rangle$\\
$S$: $\forall x, y \in H,\ \langle Tx, Ty \rangle = \langle T^* x, T^* y \rangle$
\optsfourA{$P$ and $Q$}{$P$ and $R$}{$Q$ and $S$}{$P$ and $S$}
\ans{(A)}

\textbf{Q.6} Let $X = \{a, b, c\}$ and let $\mathcal{T} = \{\phi, \{a\}, \{b\}, \{a, b\}, X\}$ be a topology defined on $X$. Then which of the following statements are TRUE?\\
$P$: $(X, \mathcal{T})$ is a Hausdorff space. \qquad $Q$: $(X, \mathcal{T})$ is a regular space.\\
$R$: $(X, \mathcal{T})$ is a normal space. \qquad $S$: $(X, \mathcal{T})$ is a connected space.
\optsfourA{$P$ and $Q$}{$Q$ and $R$}{$R$ and $S$}{$P$ and $S$}
\ans{(C)}

\textbf{Q.7} Consider the statements:\\
$P$: If $X$ is a normed linear space and $M \subseteq X$ is a subspace, then the closure $\overline{M}$ is also a subspace of $X$.\\
$Q$: If $X$ is a Banach space and $\sum x_n$ is an absolutely convergent series in $X$, then $\sum x_n$ is convergent.\\
$R$: Let $M_1$ and $M_2$ be subspaces of an inner product space such that $M_1 \cap M_2 = \{0\}$. Then $\forall m_1 \in M_1,\ m_2 \in M_2;\ \|m_1 + m_2\|^2 = \|m_1\|^2 + \|m_2\|^2$.\\
$S$: Let $f : X \to Y$ be a linear transformation from the Banach space $X$ into the Banach space $Y$. If $f$ is continuous, then the graph of $f$ is always compact.\\
The correct statements amongst the above are:
\optsfourA{$P$ and $R$ only}{$Q$ and $R$ only}{$P$ and $Q$ only}{$R$ and $S$ only}
\ans{(C)}

\textbf{Q.8} A continuous random variable $X$ has the probability density function
\[
  f(x) = \begin{cases} \dfrac{3}{5}e^{-\frac{3}{5}x}, & x > 0, \\ 0, & x \le 0. \end{cases}
\]
The probability density function of $Y = 3X + 2$ is
\optsfourA{$f(y) = \begin{cases} \dfrac{1}{5}e^{-\frac{1}{5}(y-2)}, & y > 2 \\ 0, & y \le 2 \end{cases}$}{$f(y) = \begin{cases} \dfrac{2}{5}e^{-\frac{2}{5}(y-2)}, & y > 2 \\ 0, & y \le 2 \end{cases}$}{$f(y) = \begin{cases} \dfrac{3}{5}e^{-\frac{3}{5}(y-2)}, & y > 2 \\ 0, & y \le 2 \end{cases}$}{$f(y) = \begin{cases} \dfrac{4}{5}e^{-\frac{4}{5}(y-2)}, & y > 2 \\ 0, & y \le 2 \end{cases}$}
\ans{(A)}

\textbf{Q.9} A simple random sample of size $10$ from $N(\mu, \sigma^2)$ gives $98\%$ confidence interval $(20.49, 23.51)$. Then the null hypothesis $H_0 : \mu = 20.5$ against $H_A : \mu \ne 20.5$
\optslistA{can be rejected at $2\%$ level of significance}{cannot be rejected at $5\%$ level of significance}{can be rejected at $10\%$ level of significance}{cannot be rejected at any level of significance}
\ans{(C)}

\textbf{Q.10} For the linear programming problem:
\begin{align*}
  \text{Maximize } & z = x_1 + 2x_2 + 3x_3 - 4x_4 \\
  \text{Subject to } & 2x_1 + 3x_2 - x_3 - x_4 = 15, \\
                     & 6x_1 + x_2 + x_3 - 3x_4 = 21, \\
                     & 8x_1 + 2x_2 + 3x_3 - 4x_4 = 30, \\
                     & x_1, x_2, x_3, x_4 \ge 0,
\end{align*}
$x_1 = 4$, $x_2 = 3$, $x_3 = 0$, $x_4 = 2$ is
\optslistA{an optimal solution}{a degenerate basic feasible solution}{a non-degenerate basic feasible solution}{a non-basic feasible solution}
\ans{(D)}

\textbf{Q.11} Which one of the following statements is TRUE?
\optslistA{A convex set cannot have infinitely many extreme points.}{A linear programming problem can have infinitely many extreme points.}{A linear programming problem can have exactly two different optimal solutions.}{A linear programming problem can have a non-basic optimal solution.}
\ans{(D)}

\textbf{Q.12} Let $\alpha = e^{2\pi i/5}$ and the matrix
\[
  M = \begin{bmatrix}
    1 & \alpha & \alpha^2 & \alpha^3 & \alpha^4 \\
    0 & \alpha & \alpha^2 & \alpha^3 & \alpha^4 \\
    0 & 0 & \alpha^2 & \alpha^3 & \alpha^4 \\
    0 & 0 & 0 & \alpha^3 & \alpha^4 \\
    0 & 0 & 0 & 0 & \alpha^4
  \end{bmatrix}.
\]
Then the trace of the matrix $I + M + M^2$ is
\optsfourA{$-5$}{$0$}{$3$}{$5$}
\ans{(D)}

\textbf{Q.13} Let $V = \mathbb{C}^2$ be the vector space over the field of complex numbers and $B = \{(1, i), (i, 1)\}$ be a given ordered basis of $V$. Then for which of the following is $B^* = \{f_1, f_2\}$ a dual basis of $B$ over $\mathbb{C}$?
\optslistA{$f_1(z_1, z_2) = \dfrac{1}{2}(z_1 - i z_2),\ f_2(z_1, z_2) = \dfrac{1}{2}(z_1 + i z_2)$}{$f_1(z_1, z_2) = \dfrac{1}{2}(z_1 + i z_2),\ f_2(z_1, z_2) = \dfrac{1}{2}(i z_1 + z_2)$}{$f_1(z_1, z_2) = \dfrac{1}{2}(z_1 - i z_2),\ f_2(z_1, z_2) = \dfrac{1}{2}(-i z_1 + z_2)$}{$f_1(z_1, z_2) = \dfrac{1}{2}(z_1 + i z_2),\ f_2(z_1, z_2) = \dfrac{1}{2}(-i z_1 - z_2)$}
\ans{(C)}

\textbf{Q.14} Let $R = \mathbb{Z} \times \mathbb{Z} \times \mathbb{Z}$ and $I = \mathbb{Z} \times \mathbb{Z} \times \{0\}$. Then which of the following statements is correct?
\optslistA{$I$ is a maximal ideal but not a prime ideal of $R$.}{$I$ is a prime ideal but not a maximal ideal of $R$.}{$I$ is both maximal ideal as well as a prime ideal of $R$.}{$I$ is neither a maximal ideal nor a prime ideal of $R$.}
\ans{(B)}

\textbf{Q.15} The function $u(r, \theta)$ satisfying the Laplace equation $\dfrac{\partial^2 u}{\partial r^2} + \dfrac{1}{r}\dfrac{\partial u}{\partial r} + \dfrac{1}{r^2}\dfrac{\partial^2 u}{\partial \theta^2} = 0$, $e < r < e^2$, subject to the conditions $u(e, \theta) = 1$, $u(e^2, \theta) = 0$ is
\optsfourA{$\ln(e/r)$}{$\ln(e/r^2)$}{$\ln(e^2/r)$}{$\displaystyle\sum_{n=1}^{\infty} \left(\frac{r - e^2}{e - e^2}\right)\sin n\theta$}
\ans{(C)}

\textbf{Q.16} The functional $\displaystyle\int_0^1 \left((y')^2 + (y + 2y')y'' + k x y y' + y^2\right)dx$, $y(0) = 0$, $y(1) = 1$, $y'(0) = 2$, $y'(1) = 3$, is path independent if $k$ equals
\optsfourA{$1$}{$2$}{$3$}{$4$}
\ans{(B)}

\textbf{Q.17} If a transformation $y = uv$ transforms the given differential equation $f(x)y'' - 4f'(x)y' + g(x)y = 0$ into the equation of the form $v'' + h(x)v = 0$, then $u$ must be
\optsfourA{$1/f^2$}{$xf$}{$1/2f$}{$f^2$}
\ans{(D)}

\textbf{Q.18} The expression $\dfrac{1}{D_x^2 - D_y^2}\sin(x - y)$ is equal to
\optsfourA{$-\dfrac{x}{2}\cos(x - y)$}{$-\dfrac{x}{2}\sin(x - y) + \cos(x - y)$}{$-\dfrac{x}{2}\cos(x - y) + \sin(x - y)$}{$\dfrac{3x}{2}\sin(x - y)$}
\ans{(A)}

\textbf{Q.19} The function $\phi(x)$ satisfying the integral equation $\displaystyle\int_0^x e^{x - \xi}\phi(\xi)\,d\xi = \frac{x^2}{2}$ is
\optsfourA{$\dfrac{x^2}{2}$}{$x + \dfrac{x^2}{2}$}{$x - \dfrac{x^2}{2}$}{$1 + \dfrac{x^2}{2}$}
\ans{(C)}

\textbf{Q.20} Given the data:
\begin{center}
\begin{tabular}{cccccc}
  \toprule
  $x$ & 1 & 2 & 3 & 4 & 5 \\
  \midrule
  $y$ & $-1$ & 2 & $-3$ & 4 & $-5$ \\
  \bottomrule
\end{tabular}
\end{center}
If the derivative of $y(x)$ is approximated as $y'(x_k) \approx \dfrac{1}{h}\!\left(\Delta y_k + \dfrac{1}{2}\Delta^2 y_k - \dfrac{1}{4}\Delta^3 y_k\right)$, then the value of $y'(2)$ is
\optsfourA{$4$}{$8$}{$12$}{$16$}
\ans{(B)}

\textbf{Q.21} If $A = \begin{bmatrix} 1 & 0 & 0 \\ 1 & 0 & 1 \\ 0 & 1 & 0 \end{bmatrix}$, then $A^{50}$ is
\optsfourA{$\begin{bmatrix} 1 & 0 & 0 \\ 50 & 1 & 0 \\ 50 & 0 & 1 \end{bmatrix}$}{$\begin{bmatrix} 1 & 0 & 0 \\ 48 & 1 & 0 \\ 48 & 0 & 1 \end{bmatrix}$}{$\begin{bmatrix} 1 & 0 & 0 \\ 25 & 1 & 0 \\ 25 & 0 & 1 \end{bmatrix}$}{$\begin{bmatrix} 1 & 0 & 0 \\ 24 & 1 & 0 \\ 24 & 0 & 1 \end{bmatrix}$}
\ans{(C)}

\textbf{Q.22} If $y = \displaystyle\sum_{m=0}^{\infty} c_m x^{r + m}$ is assumed to be a solution of the differential equation $x^2 y'' - xy' - 3(1 + x^2)y = 0$, then the values of $r$ are
\optsfourA{$1$ and $3$}{$-1$ and $3$}{$1$ and $-3$}{$-1$ and $-3$}
\ans{(B)}

\textbf{Q.23} Let the linear transformation $T : F^2 \to F^3$ be defined by $T(x_1, x_2) = (x_1,\, x_1 + x_2,\, x_2)$. Then the nullity of $T$ is
\optsfourA{$0$}{$1$}{$2$}{$3$}
\ans{(A)}

\textbf{Q.24} The approximate eigenvalue of the matrix
\[
  A = \begin{bmatrix}
    -15 & 4 & 3 \\
    10 & -12 & 6 \\
    20 & -4 & 2
  \end{bmatrix}
\]
obtained after two iterations of Power method, with the initial vector $[1\;\;1\;\;1]^T$, is
\optsfourA{$7.768$}{$9.468$}{$10.548$}{$19.468$}
\ans{(C)}

\textbf{Q.25} The root of the equation $x e^x = 1$ between $0$ and $1$, obtained by using two iterations of bisection method, is
\optsfourA{$0.25$}{$0.50$}{$0.75$}{$0.65$}
\ans{(C)}

\section*{Q.26 -- Q.47 Carry Two Marks Each}

\textbf{Q.26} Let $\displaystyle\oint_C \left[\frac{1}{(z - 2)^2} - \frac{a - 2}{(z - 2)} + b\right]dz = 4\pi i$, where the closed curve $C$ is the triangle having vertices at $i$, $\dfrac{-1 - i}{\sqrt{2}}$, and $\dfrac{1 - i}{\sqrt{2}}$, the integral being taken in anti-clockwise direction. Then one value of $a$ is
% NOTE: Page 191 in source PDF was mirrored; verify exact integrand and vertices of Q.26 against original paper.
\optsfourA{$1 + i$}{$2 + i$}{$3 + i$}{$4 + i$}
\ans{(C)}

\textbf{Q.27} The Lebesgue measure of the set $A = \left\{0 < x < 1 : x\sin\!\left(\dfrac{\pi}{2x}\right) \ge 0\right\}$ is
% NOTE: Page 191 in source PDF was mirrored; verify exact argument of sin against original paper.
\optsfourA{$0$}{$1$}{$\ln 2$}{$1 - \ln\sqrt{2}$}
\ans{(D)}

\textbf{Q.28} Which of the following statements are TRUE?\\
$P$: The set $\left\{x \in \mathbb{R} : |\cos x| \le \dfrac{1}{2}\right\}$ is compact.\\
$Q$: The set $\{x \in \mathbb{R} : \tan x \text{ is not differentiable}\}$ is complete.\\
$R$: The set $\left\{x \in \mathbb{R} : \displaystyle\sum_{n=0}^{\infty} \frac{(-1)^n x^{2n+1}}{(2n+1)!} \text{ is convergent}\right\}$ is bounded.\\
$S$: The set $\{x \in \mathbb{R} : f(x) = \cos x \text{ has a local maxima}\}$ is closed.
\optsfourA{$P$ and $Q$}{$R$ and $S$}{$Q$ and $S$}{$P$ and $S$}
\ans{(C)}

\textbf{Q.29} If a random variable $X$ assumes only positive integral values, with the probability $P(X = x) = \dfrac{2}{3}\left(\dfrac{1}{3}\right)^{x-1}$, $x = 1, 2, 3, \dots$, then $E(X)$ is
\optsfourA{$2/9$}{$2/3$}{$1$}{$3/2$}
\ans{(D)}

\textbf{Q.30} The probability density function of the random variable $X$ is
\[
  f(x) = \begin{cases} \dfrac{1}{\lambda}e^{-x/\lambda}, & x > 0, \\ 0, & x \le 0, \end{cases}
\]
where $\lambda > 0$. For testing the hypothesis $H_0 : \lambda = 3$ against $H_A : \lambda = 5$, a test is given as ``Reject $H_0$ if $X \ge 4.5$''. The probability of type I error and power of this test are, respectively,
\optsfourA{$0.1353$ and $0.4966$}{$0.1827$ and $0.379$}{$0.2021$ and $0.4493$}{$0.2231$ and $0.4066$}
\ans{(D)}

\textbf{Q.31} The order of the smallest possible non-trivial group containing elements $x$ and $y$ such that $x^7 = y^2 = e$ and $yx = x^4 y$ is
\optsfourA{$1$}{$2$}{$7$}{$14$}
\ans{(B)}

\textbf{Q.32} The number of $5$-Sylow subgroup(s) in a group of order $45$ is
\optsfourA{$1$}{$2$}{$3$}{$4$}
\ans{(A)}

\textbf{Q.33} The solution of the initial value problem $y'' + 2y' + 10y = 6\delta(t)$, $y(0) = 0$, $y'(0) = 0$, where $\delta(t)$ denotes the Dirac-delta function, is
\optsfourA{$2e^t\sin 3t$}{$6e^t\sin 3t$}{$2e^{-t}\sin 3t$}{$6e^{-t}\sin 3t$}
\ans{Marks to All (official key; mathematically (C) $2e^{-t}\sin 3t$ for $t > 0$)}

\textbf{Q.34} Let $\omega = \cos\dfrac{2\pi}{3} + i\sin\dfrac{2\pi}{3}$, $M = \begin{pmatrix} 0 & i \\ i & 0 \end{pmatrix}$, $N = \begin{pmatrix} \omega & 0 \\ 0 & \omega^2 \end{pmatrix}$, and $G = \langle M, N \rangle$ be the group generated by the matrices $M$ and $N$ under matrix multiplication. Then
\optsfourA{$G / Z(G) \cong C_6$}{$G / Z(G) \cong S_3$}{$G / Z(G) \cong C_2$}{$G / Z(G) \cong C_4$}
\ans{(B)}

\textbf{Q.35} The flux of the vector field $\vec{u} = x\hat{i} + y\hat{j} + z\hat{k}$ flowing out through the surface of the ellipsoid $\dfrac{x^2}{a^2} + \dfrac{y^2}{b^2} + \dfrac{z^2}{c^2} = 1$, $a > b > c > 0$, is
\optsfourA{$\pi abc$}{$2\pi abc$}{$3\pi abc$}{$4\pi abc$}
\ans{(D)}

\textbf{Q.36} The integral surface satisfying the partial differential equation $\dfrac{\partial z}{\partial x} + z^2\dfrac{\partial z}{\partial y} = 0$ and passing through the straight line $x = 1$, $y = z$ is
\optsfourA{$(x - 1)z + z^2 = y^2$}{$x^2 + y^2 - z^2 = 1$}{$(y - z)x + x^2 = 1$}{$(x - 1)z^2 + z = y$}
\ans{(D)}

\textbf{Q.37} The diffusion equation $\dfrac{\partial^2 u}{\partial x^2} = \dfrac{\partial u}{\partial t}$, $u = u(x, t)$, $u(0, t) = 0 = u(\pi, t)$, $u(x, 0) = \cos x \sin 5x$, admits the solution
\optsfourA{$\dfrac{e^{-36t}}{2}[\sin 6x + e^{20t}\sin 4x]$}{$\dfrac{e^{-36t}}{2}[\sin 4x + e^{20t}\sin 6x]$}{$\dfrac{e^{-20t}}{2}[\sin 3x + e^{15t}\sin 5x]$}{$\dfrac{e^{-36t}}{2}[\sin 5x + e^{20t}\sin x]$}
\ans{(A)}

\textbf{Q.38} Let $f(x)$ and $x f(x)$ be the particular solutions of a differential equation $y'' + R(x)y' + S(x)y = 0$. Then the solution of the differential equation $y'' + R(x)y' + S(x)y = f(x)$ is
\optsfourA{$y = \left(-\dfrac{x^2}{2} + \alpha x + \beta\right)f(x)$}{$y = \left(\dfrac{x^2}{2} + \alpha x + \beta\right)f(x)$}{$y = (-x^2 + \alpha x + \beta)f(x)$}{$y = (x^3 + \alpha x + \beta)f(x)$}
\ans{(B)}

\textbf{Q.39} Let the Legendre equation $(1 - x^2)y'' - 2xy' + n(n + 1)y = 0$ have $n$-th degree polynomial solution $y_n(x)$ such that $y_n(1) = 3$. If $\displaystyle\int_{-1}^{1} \left(y_n^2(x) + y_{n-1}^2(x)\right)dx = \frac{144}{15}$, then $n$ is
\optsfourA{$1$}{$2$}{$3$}{$4$}
\ans{(B)}

\textbf{Q.40} The maximum value of the function $f(x, y, z) = xyz$ subject to the constraint $xy + yz + zx - a = 0$, $a > 0$, is
\optsfourA{$a^{3/2}$}{$(a/3)^{3/2}$}{$(3/a)^{3/2}$}{$(3a/2)^{3/2}$}
\ans{(B)}

\textbf{Q.41} The functional $\displaystyle\int_0^1 \left((y')^2 + 4y^2 + 8y e^x\right)dx$, $y(0) = -\dfrac{4}{3}$, $y(1) = -\dfrac{4e}{3}$, possesses:
\optsfourA{strong minima on $y = -\dfrac{1}{3}e^x$}{strong minima on $y = -\dfrac{4}{3}e^x$}{weak maxima on $y = -\dfrac{1}{3}e^x$}{strong maxima on $y = -\dfrac{4}{3}e^x$}
\ans{(B)}

\textbf{Q.42} A particle of mass $m$ is constrained to move on a circle with radius $a$ which itself is rotating about its vertical diameter with a constant angular velocity $\omega$. Assume that the initial angular velocity is zero and $g$ is the acceleration due to gravity. If $\theta$ be the inclination of the radius vector of the particle with the axis of rotation and $\dot{\theta}$ denotes the derivative of $\theta$ with respect to $t$, then the Lagrangian of this system is
\optslistA{$\dfrac{1}{2}ma^2(\dot{\theta}^2 + \omega^2\sin^2\theta) + mga\cos\theta$}{$\dfrac{1}{2}ma^2(\dot{\theta}^2 + 2\omega\sin\theta) - mga\sin\theta$}{$\dfrac{1}{2}ma^2(\dot{\theta}^2 + 2\omega^2\cos\theta) - mga\sin\theta$}{$\dfrac{1}{2}ma^2(\dot{\theta}^2 + \omega\sin 2\theta) + mga\sin\theta$}
\ans{(A)}

\textbf{Q.43} For the matrix
\[
  M = \begin{bmatrix} 2 & 3 + 2i & -4 \\ 3 - 2i & 5 & 6i \\ -4 & -6i & 3 \end{bmatrix},
\]
which of the following statements are correct?\\
$P$: $M$ is skew-Hermitian and $iM$ is Hermitian.\\
$Q$: $M$ is Hermitian and $iM$ is skew-Hermitian.\\
$R$: eigenvalues of $M$ are real.\\
$S$: eigenvalues of $iM$ are real.
\optsfourA{$P$ and $R$ only}{$Q$ and $R$ only}{$P$ and $S$ only}{$Q$ and $S$ only}
\ans{(B)}

\textbf{Q.44} Let $T : P_3 \to P_3$ be the map given by $T(p(x)) = \displaystyle\int_1^x p'(t)\,dt$. If the matrix of $T$ relative to the standard bases $B_1 = B_2 = \{1, x, x^2, x^3\}$ is $M$ and $M'$ the transpose of the matrix $M$, then $M + M'$ is
\optsfourA{$\begin{bmatrix} 0 & -1 & -1 & -1 \\ -1 & 2 & 0 & 0 \\ 0 & 0 & 2 & 0 \\ 0 & 0 & 0 & 2 \end{bmatrix}$}{$\begin{bmatrix} -1 & 0 & 0 & 2 \\ 0 & -1 & 1 & 0 \\ 0 & 1 & -1 & 0 \\ 2 & 0 & 2 & -1 \end{bmatrix}$}{$\begin{bmatrix} 2 & 0 & 0 & -1 \\ 0 & 2 & 1 & 0 \\ 0 & 1 & 2 & -1 \\ -1 & 0 & -1 & 0 \end{bmatrix}$}{$\begin{bmatrix} 0 & 2 & 2 & 2 \\ 2 & -1 & 0 & 0 \\ 2 & 0 & -1 & 0 \\ 2 & 0 & 0 & -1 \end{bmatrix}$}
\ans{Marks to All (official key)}

\textbf{Q.45} Using Euler's method taking step size $= 0.1$, the approximate value of $y$ obtained corresponding to $x = 0.2$ for the initial value problem $\dfrac{dy}{dx} = x^2 + y^2$ and $y(0) = 1$, is
\optsfourA{$1.322$}{$1.122$}{$1.222$}{$1.110$}
\ans{(C)}

\textbf{Q.46} The following table gives the unit transportation costs, the supply at each origin and the demand of each destination for a transportation problem:
\begin{center}
\begin{tabular}{lccccc}
  \toprule
  Origin & $D_1$ & $D_2$ & $D_3$ & $D_4$ & Supply \\
  \midrule
  $O_1$ & 3 & 4 & 8 & 7 & 60 \\
  $O_2$ & 7 & 3 & 7 & 6 & 80 \\
  $O_3$ & 3 & 9 & 3 & 4 & 100 \\
  \midrule
  Demand & 40 & 70 & 50 & 80 & \\
  \bottomrule
\end{tabular}
\end{center}
Let $x_{ij}$ denote the number of units to be transported from origin $i$ to destination $j$. If the $u$--$v$ method is applied to improve the basic feasible solution given by $x_{12} = 60$, $x_{22} = 10$, $x_{23} = 50$, $x_{24} = 20$, $x_{31} = 40$, and $x_{34} = 60$, then the variables entering and leaving the basis, respectively, are
\optsfourA{$x_{11}$ and $x_{24}$}{$x_{13}$ and $x_{23}$}{$x_{14}$ and $x_{24}$}{$x_{33}$ and $x_{24}$}
\ans{(A)}

\textbf{Q.47} Consider the system of equations
\[
  \begin{bmatrix} 5 & -1 & 1 \\ 2 & 4 & 0 \\ 1 & 1 & 5 \end{bmatrix}
  \begin{bmatrix} x \\ y \\ z \end{bmatrix}
  =
  \begin{bmatrix} 10 \\ 12 \\ -1 \end{bmatrix}.
\]
Using Jacobi's method with the initial guess $[x^{(0)}\;\;y^{(0)}\;\;z^{(0)}]^T = [2.0\;\;3.0\;\;0.0]^T$, the approximate solution $[x^{(2)}\;\;y^{(2)}\;\;z^{(2)}]^T$ after two iterations is
\optsfourA{$[2.64\;\;-1.70\;\;-1.12]^T$}{$[2.64\;\;-1.70\;\;1.12]^T$}{$[2.64\;\;1.70\;\;-1.12]^T$}{$[2.64\;\;1.70\;\;1.12]^T$}
\ans{(C)}

\section*{Common Data Questions}

\textbf{Common Data for Questions 48 and 49:}
The optimal table for the primal linear programming problem: Maximize $z = 6x_1 + 12x_2 + 12x_3 - 6x_4$ subject to $x_1 + x_2 + x_3 = 4$, $x_1 + 4x_2 + x_4 = 8$, $x_1, x_2, x_3, x_4 \ge 0$, is
\begin{center}
\begin{tabular}{cccccc}
  \toprule
  Basic variables ($x_B$) & $x_1$ & $x_2$ & $x_3$ & $x_4$ & RHS ($b$) \\
  \midrule
  $x_3$ & $3/4$ & $0$ & $1$ & $-1/4$ & $2$ \\
  $x_2$ & $1/4$ & $1$ & $0$ & $1/4$ & $2$ \\
  \midrule
  $z_j - c_j$ & $6$ & $0$ & $0$ & $6$ & $z = 48$ \\
  \bottomrule
\end{tabular}
\end{center}

\textbf{Q.48} If $y_1$ and $y_2$ are the dual variables corresponding to the first and second primal constraints, then their values in the optimal solution of the dual problem are, respectively,
\optsfourA{$0$ and $6$}{$12$ and $0$}{$6$ and $3$}{$4$ and $4$}
\ans{(B)}

\textbf{Q.49} If the right hand side of the second constraint is changed from $8$ to $20$, then in the optimal solution of the primal problem, the basic variables will be
\optsfourA{$x_1$ and $x_2$}{$x_1$ and $x_3$}{$x_2$ and $x_3$}{$x_2$ and $x_4$}
\ans{(D)}

\medskip
\textbf{Common Data for Questions 50 and 51:}
Consider the Fredholm integral equation $u(x) = x + \lambda \displaystyle\int_0^1 x e^t u(t)\,dt$.

\textbf{Q.50} The resolvent kernel $R(x, t; \lambda)$ for this integral equation is
\optsfourA{$\dfrac{x e^t}{1 - \lambda}$}{$\dfrac{\lambda x e^t}{1 + \lambda}$}{$\dfrac{x e^t}{1 + \lambda^2}$}{$\dfrac{x e^t}{1 - \lambda^2}$}
\ans{(A)}

\textbf{Q.51} The solution of this integral equation is
\optsfourA{$\dfrac{x + 1}{1 - \lambda}$}{$\dfrac{x^2}{1 - \lambda^2}$}{$\dfrac{x}{1 + \lambda^2}$}{$\dfrac{x}{1 - \lambda}$}
\ans{(D)}

\section*{Linked Answer Questions}

\textbf{Statement for Linked Answer Questions 52 and 53:}
The joint probability density function of two random variables $X$ and $Y$ is given as
\[
  f(x, y) = \begin{cases} \dfrac{6}{5}(x + y^2), & 0 \le x \le 1,\ 0 \le y \le 1, \\ 0, & \text{elsewhere.} \end{cases}
\]

\textbf{Q.52} $E(X)$ and $E(Y)$ are, respectively,
\optsfourA{$\dfrac{2}{5}$ and $\dfrac{3}{5}$}{$\dfrac{3}{5}$ and $\dfrac{3}{5}$}{$\dfrac{3}{5}$ and $\dfrac{6}{5}$}{$\dfrac{4}{5}$ and $\dfrac{6}{5}$}
\ans{(B)}

\textbf{Q.53} $\operatorname{Cov}(X, Y)$ is
\optsfourA{$-0.01$}{$0$}{$0.01$}{$0.02$}
\ans{(A)}

\medskip
\textbf{Statement for Linked Answer Questions 54 and 55:}
Consider the functions $f(z) = \dfrac{z^2 + \alpha z}{(z + 1)^2}$ and $g(z) = \sinh\!\left(z - \dfrac{\pi}{2\alpha}\right)$, $\alpha \ne 0$.

\textbf{Q.54} The residue of $f(z)$ at its pole is equal to $1$. Then the value of $\alpha$ is
\optsfourA{$-1$}{$1$}{$2$}{$3$}
\ans{(D)}

\textbf{Q.55} For the value of $\alpha$ obtained in Q.54, the function $g(z)$ is not conformal at a point
\optsfourA{$\dfrac{\pi(1 + 3i)}{6}$}{$\dfrac{\pi(3 + i)}{6}$}{$\dfrac{2\pi}{3}$}{$\dfrac{i\pi}{2}$}
\ans{(A)}

\section*{General Aptitude (GA) Questions}

\subsection*{Q.56 -- Q.60 Carry One Mark Each}

\textbf{Q.56} Choose the most appropriate word from the options given below to complete the following sentence:\\
Given the seriousness of the situation that he had to face, his \underline{\hspace{2cm}} was impressive.
\optsfourA{beggary}{nomenclature}{jealousy}{nonchalance}
\ans{(D)}

\textbf{Q.57} Choose the most appropriate alternative from the options given below to complete the following sentence:\\
If the tired soldier wanted to lie down, he \underline{\hspace{2cm}} the mattress out on the balcony.
\optsfourA{should take}{shall take}{should have taken}{will have taken}
\ans{(A)}

\textbf{Q.58} If $(1.001)^{1259} = 3.52$ and $(1.001)^{2062} = 7.85$, then $(1.001)^{3321} =$
\optsfourA{$2.23$}{$4.33$}{$11.37$}{$27.64$}
\ans{(D)}

\textbf{Q.59} One of the parts (A, B, C, D) in the sentence given below contains an ERROR. Which one of the following is INCORRECT?\\
\textit{I requested that he should be given the driving test today instead of tomorrow.}
\optsfourA{requested that}{should be given}{the driving test}{instead of tomorrow}
\ans{(B)}

\textbf{Q.60} Which one of the following options is the closest in meaning to the word given below?\\
\textbf{Latitude}
\optsfourA{Eligibility}{Freedom}{Coercion}{Meticulousness}
\ans{(B)}

\subsection*{Q.61 -- Q.65 Carry Two Marks Each}

\textbf{Q.61} There are eight bags of rice looking alike, seven of which have equal weight and one is slightly heavier. The weighing balance is of unlimited capacity. Using this balance, the minimum number of weighings required to identify the heavier bag is
\optsfourA{$2$}{$3$}{$4$}{$8$}
\ans{(A)}

\textbf{Q.62} Raju has $14$ currency notes in his pocket consisting of only Rs.\ $20$ notes and Rs.\ $10$ notes. The total money value of the notes is Rs.\ $230$. The number of Rs.\ $10$ notes that Raju has is
\optsfourA{$5$}{$6$}{$9$}{$10$}
\ans{(A)}

\textbf{Q.63} One of the legacies of the Roman legions was discipline. In the legions, military law prevailed and discipline was brutal. Discipline on the battlefield kept units obedient, intact and fighting, even when the odds and conditions were against them.\\
Which one of the following statements best sums up the meaning of the above passage?
\optslistA{Thorough regimentation was the main reason for the efficiency of the Roman legions even in adverse circumstances.}{The legions were treated inhumanly as if the men were animals.}{Discipline was the armies' inheritance from their seniors.}{The harsh discipline to which the legions were subjected to led to the odds and conditions being against them.}
\ans{(A)}

\textbf{Q.64} $A$ and $B$ are friends. They decide to meet between $1$ PM and $2$ PM on a given day. There is a condition that whoever arrives first will not wait for the other for more than $15$ minutes. The probability that they will meet on that day is
\optsfourA{$1/4$}{$1/16$}{$7/16$}{$9/16$}
\ans{(C)}

\textbf{Q.65} The data given in the following table summarizes the monthly budget of an average household:
\begin{center}
\begin{tabular}{lc}
  \toprule
  Category & Amount (Rs.) \\
  \midrule
  Food & 4000 \\
  Clothing & 1200 \\
  Rent & 2000 \\
  Savings & 1500 \\
  Other expenses & 1800 \\
  \bottomrule
\end{tabular}
\end{center}
The approximate percentage of the monthly budget NOT spent on savings is
\optsfourA{$10\%$}{$14\%$}{$81\%$}{$86\%$}
\ans{(D)}

\newpage

% ============================================================================
% GATE 2013 -- MATHEMATICS (MA) & GENERAL APTITUDE (GA)
% ============================================================================

\begin{center}
  {\Large\bfseries GATE 2013 -- Mathematics (MA)}\\[4pt]
  {\small Duration: Three Hours \quad\textbar\quad Maximum Marks: 100}
\end{center}

\medskip
\textbf{Notations:} $\mathbb{R}, \mathbb{Z}, \mathbb{C}, \mathbb{N}, \mathbb{Q}$: standard sets; $\mathbb{Z}_n$: cyclic group of order $n$; $S_n$: permutation group on $\{1,\dots,n\}$; $X^t$: transpose of $X$; $\operatorname{Re}(z), \operatorname{Im}(z)$: real and imaginary parts; $\overline{A}, A^\circ$: closure and interior; $\langle a \rangle$: ideal generated by $a$; $P(E), E[X], \operatorname{Var}(X)$: probability, expectation, variance; $\log x$: natural logarithm.

\section*{Q.1 -- Q.25 Carry One Mark Each}

\textbf{Q.1} The possible set of eigenvalues of a $4 \times 4$ skew-symmetric orthogonal real matrix is
\optsfourA{$\{\pm i\}$}{$\{\pm i, \pm 1\}$}{$\{\pm 1\}$}{$\{0, \pm i\}$}
\ans{(A)}

\textbf{Q.2} The coefficient of $(z - \pi)^2$ in the Taylor series expansion of
\[
  f(z) = \begin{cases} \dfrac{\sin z}{z - \pi}, & \text{if } z \ne \pi, \\ -1, & \text{if } z = \pi, \end{cases}
\]
around $\pi$ is
\optsfourA{$\dfrac{1}{2}$}{$-\dfrac{1}{2}$}{$\dfrac{1}{6}$}{$-\dfrac{1}{6}$}
\ans{(C)}

\textbf{Q.3} Consider $\mathbb{R}^2$ with the usual topology. Which of the following statements are TRUE for all $A, B \subseteq \mathbb{R}^2$?\\
$P$: $\overline{A \cup B} = \overline{A} \cup \overline{B}$. \qquad $Q$: $\overline{A \cap B} = \overline{A} \cap \overline{B}$.\\
$R$: $(A \cup B)^\circ = A^\circ \cup B^\circ$. \qquad $S$: $(A \cap B)^\circ = A^\circ \cap B^\circ$.
\optsfourA{$P$ and $R$ only}{$P$ and $S$ only}{$Q$ and $R$ only}{$Q$ and $S$ only}
\ans{(B)}

\textbf{Q.4} Let $f : \mathbb{R} \to \mathbb{R}$ be a continuous function with $f(1) = 5$ and $f(3) = 11$. If $g(x) = \displaystyle\int_1^3 f(x + t)\,dt$, then $g'(0)$ is equal to \underline{\hspace{2cm}}.
\ans{$6$}

\textbf{Q.5} Let $P$ be a $2 \times 2$ complex matrix such that $\operatorname{trace}(P) = 1$ and $\det(P) = -6$. Then $\operatorname{trace}(P^4 - P^3)$ is \underline{\hspace{2cm}}.
\ans{$78$}

\textbf{Q.6} Suppose that $R$ is a unique factorization domain and that $a, b \in R$ are distinct irreducible elements. Which of the following statements is TRUE?
\optslistA{The ideal $\langle 1 + a \rangle$ is a prime ideal}{The ideal $\langle a + b \rangle$ is a prime ideal}{The ideal $\langle 1 + ab \rangle$ is a prime ideal}{The ideal $\langle a \rangle$ is not necessarily a maximal ideal}
\ans{(D)}

\textbf{Q.7} Let $X$ be a compact Hausdorff topological space and let $Y$ be a topological space. Let $f : X \to Y$ be a bijective continuous mapping. Which of the following is TRUE?
\optslistA{$f$ is a closed map but not necessarily an open map}{$f$ is an open map but not necessarily a closed map}{$f$ is both an open map and a closed map}{$f$ need not be an open map or a closed map}
\ans{(D)}

\textbf{Q.8} Consider the linear programming problem: Maximize $x + \dfrac{3}{2}y$ subject to $2x + 3y \le 16$, $x + 4y \le 18$, $x \ge 0$, $y \ge 0$. If $S$ denotes the set of all solutions of the above problem, then
\optsfourA{$S$ is empty}{$S$ is a singleton}{$S$ is a line segment}{$S$ has positive area}
\ans{(C)}

\textbf{Q.9} Which of the following groups has a proper subgroup that is NOT cyclic?
\optsfourA{$\mathbb{Z}_{15} \times \mathbb{Z}_{77}$}{$S_3$}{$(\mathbb{Z}, +)$}{$(\mathbb{Q}, +)$}
\ans{(D)}

\textbf{Q.10} The value of the integral $\displaystyle\int_0^\infty \int_x^\infty \left(\frac{1}{y}\right) e^{-y/2}\,dy\,dx$ is \underline{\hspace{2cm}}.
\ans{$2$}

\textbf{Q.11} Suppose the random variable $U$ has uniform distribution on $[0, 1]$ and $X = -2\log U$. The density of $X$ is
\optsfourA{$f(x) = \begin{cases} e^{-x}, & x > 0 \\ 0, & \text{otherwise} \end{cases}$}{$f(x) = \begin{cases} 2e^{-2x}, & x > 0 \\ 0, & \text{otherwise} \end{cases}$}{$f(x) = \begin{cases} \dfrac{1}{2}e^{-x/2}, & x > 0 \\ 0, & \text{otherwise} \end{cases}$}{$f(x) = \begin{cases} 1/2, & x \in [0, 2] \\ 0, & \text{otherwise} \end{cases}$}
\ans{(C)}

\textbf{Q.12} Let $f$ be an entire function on $\mathbb{C}$ such that $|f(z)| \le 100\log|z|$ for each $z$ with $|z| \ge 2$. If $f(i) = 2i$, then $f(1)$
\optsfourA{must be $2$}{must be $2i$}{must be $i$}{cannot be determined from the given data}
\ans{(B)}

\textbf{Q.13} The number of group homomorphisms from $\mathbb{Z}_3$ to $\mathbb{Z}_9$ is \underline{\hspace{2cm}}.
\ans{$3$}

\textbf{Q.14} Let $u(x, t)$ be the solution to the wave equation $\dfrac{\partial^2 u}{\partial x^2}(x, t) = \dfrac{\partial^2 u}{\partial t^2}(x, t)$, $u(x, 0) = \cos(5\pi x)$, $\dfrac{\partial u}{\partial t}(x, 0) = 0$. Then the value of $u(1, 1)$ is \underline{\hspace{2cm}}.
\ans{$1$}

\textbf{Q.15} Let $f(x) = \displaystyle\sum_{n=1}^{\infty} \frac{\sin(nx)}{n^2}$. Then
\optsfourA{$\lim_{x\to 0} f(x) = 0$}{$\lim_{x\to 0} f(x) = 1$}{$\lim_{x\to 0} f(x) = \pi^2/6$}{$\lim_{x\to 0} f(x)$ does not exist}
\ans{(A)}

\textbf{Q.16} Suppose $X$ is a random variable with $P(X = k) = (1 - p)^k p$ for $k \in \{0, 1, 2, \dots\}$ and some $p \in (0, 1)$. For the hypothesis testing problem $H_0 : p = \dfrac{1}{2}$ vs.\ $H_1 : p \ne \dfrac{1}{2}$, consider the test ``Reject $H_0$ if $X \le A$ or if $X \ge B$'', where $A < B$ are given positive integers. The type-I error of this test is
\optsfourA{$1 + 2^{-B} - 2^{-A}$}{$1 - 2^{-B} + 2^{-A}$}{$1 + 2^{-B} - 2^{-A-1}$}{$1 - 2^{-B} + 2^{-A-1}$}
\ans{(C)}

\textbf{Q.17} Let $G$ be a group of order $231$. The number of elements of order $11$ in $G$ is \underline{\hspace{2cm}}.
\ans{$10$}

\textbf{Q.18} Let $f : \mathbb{R}^2 \to \mathbb{R}^2$ be defined by $f(x, y) = (e^{x+y}, e^{x-y})$. The area of the image of the region $\{(x, y) \in \mathbb{R}^2 : 0 < x, y < 1\}$ under the mapping $f$ is
\optsfourA{$1$}{$e - 1$}{$e^2$}{$e^2 - 1$}
\ans{(D)}

\textbf{Q.19} Which of the following is a field?
\optsfourA{$\mathbb{C}[x]/\langle x^2 + 2 \rangle$}{$\mathbb{Z}[x]/\langle x^2 + 2 \rangle$}{$\mathbb{Q}[x]/\langle x^2 - 2 \rangle$}{$\mathbb{R}[x]/\langle x^2 - 2 \rangle$}
\ans{(C)}

\textbf{Q.20} Let $x_0 = 0$. Define $x_{n+1} = \cos x_n$ for every $n \ge 0$. Then
\optslistA{$\{x_n\}$ is increasing and convergent}{$\{x_n\}$ is decreasing and convergent}{$\{x_n\}$ is convergent and $x_{2n} < \lim_{m\to\infty} x_m < x_{2n+1}$ for every $n \in \mathbb{N}$}{$\{x_n\}$ is not convergent}
\ans{(C)}

\textbf{Q.21} Let $C$ be the contour $|z| = 2$ oriented in the anti-clockwise direction. The value of the integral $\displaystyle\oint_C z e^{3/z}\,dz$ is
\optsfourA{$3\pi i$}{$5\pi i$}{$7\pi i$}{$9\pi i$}
\ans{(D)}

\textbf{Q.22} For each $\lambda > 0$, let $X_\lambda$ be a random variable with exponential density $\lambda e^{-\lambda x}$ on $(0, \infty)$. Then $\operatorname{Var}(\log X_\lambda)$
\optsfourA{is strictly increasing in $\lambda$}{is strictly decreasing in $\lambda$}{does not depend on $\lambda$}{first increases and then decreases in $\lambda$}
\ans{(C)}

\textbf{Q.23} Let $\{a_n\}$ be the sequence of consecutive positive solutions of the equation $\tan x = x$ and let $\{b_n\}$ be the sequence of consecutive positive solutions of the equation $\tan\sqrt{x} = x$. Then
\optsfourA{$\sum_{n=1}^{\infty} \dfrac{1}{a_n}$ converges but $\sum_{n=1}^{\infty} \dfrac{1}{b_n}$ diverges}{$\sum_{n=1}^{\infty} \dfrac{1}{a_n}$ diverges but $\sum_{n=1}^{\infty} \dfrac{1}{b_n}$ converges}{Both $\sum_{n=1}^{\infty} \dfrac{1}{a_n}$ and $\sum_{n=1}^{\infty} \dfrac{1}{b_n}$ converge}{Both $\sum_{n=1}^{\infty} \dfrac{1}{a_n}$ and $\sum_{n=1}^{\infty} \dfrac{1}{b_n}$ diverge}
\ans{(B)}

\textbf{Q.24} Let $f$ be an analytic function on $\overline{D} = \{z \in \mathbb{C} : |z| \le 1\}$. Assume that $|f(z)| \le 1$ for each $z \in \overline{D}$. Then, which of the following is NOT a possible value of $(e^f)''(0)$?
\optsfourA{$2$}{$6$}{$\dfrac{7}{9}e^{1/9}$}{$\sqrt{2} + i\sqrt{2}$}
\ans{(B)}

\textbf{Q.25} The number of non-isomorphic abelian groups of order $24$ is \underline{\hspace{2cm}}.
\ans{$3$}

\section*{Q.26 -- Q.47 Carry Two Marks Each}

\textbf{Q.26} Let $V$ be the real vector space of all polynomials in one variable with real coefficients and having degree at most $20$. Define the subspaces
\[
  W_1 = \left\{p \in V : p(1) = 0,\ p\!\left(\frac{1}{2}\right) = 0,\ p(5) = 0,\ p(7) = 0\right\},
\]
\[
  W_2 = \left\{p \in V : p\!\left(\frac{1}{2}\right) = 0,\ p(3) = 0,\ p(4) = 0,\ p(7) = 0\right\}.
\]
Then the dimension of $W_1 \cap W_2$ is \underline{\hspace{2cm}}.
\ans{$15$}

\textbf{Q.27} Let $f, g : [0, 1] \to \mathbb{R}$ be defined by
\[
  f(x) = \begin{cases} x, & \text{if } x = \dfrac{1}{n} \text{ for } n \in \mathbb{N}, \\ 0, & \text{otherwise,} \end{cases}
  \qquad \text{and} \qquad
  g(x) = \begin{cases} 1, & \text{if } x \in \mathbb{Q} \cap [0, 1], \\ 0, & \text{otherwise.} \end{cases}
\]
Then
\optslistA{Both $f$ and $g$ are Riemann integrable}{$f$ is Riemann integrable and $g$ is Lebesgue integrable}{$g$ is Riemann integrable and $f$ is Lebesgue integrable}{Neither $f$ nor $g$ is Riemann integrable}
\ans{(B)}

\textbf{Q.28} Consider the following linear programming problem: Maximize $x + 3y + 6z - w$ subject to $5x + y + 6z + 7w \le 20$, $x + 2y + 2z + 9w \le 40$, $x \ge 0$, $y \ge 0$, $z \ge 0$, $w \ge 0$. Then the optimal value is \underline{\hspace{2cm}}.
\ans{$60$}

\textbf{Q.29} Suppose $X$ is a real-valued random variable. Which of the following values CANNOT be attained by $E[X]$ and $E[X^2]$, respectively?
\optsfourA{$0$ and $1$}{$2$ and $3$}{$\dfrac{1}{2}$ and $\dfrac{1}{3}$}{$2$ and $5$}
\ans{(B)}

\textbf{Q.30} Which of the following subsets of $\mathbb{R}^2$ is NOT compact?
\optslistA{$\{(x, y) \in \mathbb{R}^2 : -1 \le x \le 1,\ y = \sin x\}$}{$\{(x, y) \in \mathbb{R}^2 : -1 \le y \le 1,\ y = x^8 - x^3 - 1\}$}{$\{(x, y) \in \mathbb{R}^2 : y = 0,\ \sin(e^{-x}) = 0\}$}{$\{(x, y) \in \mathbb{R}^2 : x > 0,\ y = \sin(1/x)\} \cap \{(x, y) \in \mathbb{R}^2 : x > 0,\ y = 1/x\}$}
\ans{(C)}

\textbf{Q.31} Let $M$ be the real vector space of $2 \times 3$ matrices with real entries. Let $T : M \to M$ be defined by
\[
  T\!\left(\begin{bmatrix} x_1 & x_2 & x_3 \\ x_4 & x_5 & x_6 \end{bmatrix}\right) = \begin{bmatrix} -x_6 & x_4 & x_1 \\ x_3 & x_5 & x_2 \end{bmatrix}.
\]
The determinant of $T$ is \underline{\hspace{2cm}}.
\ans{$-1$}

\textbf{Q.32} Let $\mathcal{H}$ be a Hilbert space and let $\{e_n : n \ge 1\}$ be an orthonormal basis of $\mathcal{H}$. Suppose $T : \mathcal{H} \to \mathcal{H}$ is a bounded linear operator. Which of the following CANNOT be true?
\optslistA{$T(e_n) = e_1$ for all $n \ge 1$}{$T(e_n) = e_{n+1}$ for all $n \ge 1$}{$T(e_n) = \sqrt{\dfrac{n + 1}{n}}\,e_n$ for all $n \ge 1$}{$T(e_n) = e_{n-1}$ for all $n \ge 2$ and $T(e_1) = 0$}
\ans{(A)}

\textbf{Q.33} The value of the limit $\displaystyle\lim_{n\to\infty} \frac{2^{-n^2}}{\sum_{k=n+1}^{\infty} 2^{-k^2}}$ is
\optsfourA{$0$}{some $c \in (0, 1)$}{$1$}{$\infty$}
\ans{(D)}

\textbf{Q.34} Let $f : \mathbb{C} \setminus \{3i\} \to \mathbb{C}$ be defined by $f(z) = \dfrac{z - i}{iz + 3}$. Which of the following statements about $f$ is FALSE?
\optslistA{$f$ is conformal on $\mathbb{C} \setminus \{3i\}$}{$f$ maps circles in $\mathbb{C} \setminus \{3i\}$ onto circles in $\mathbb{C}$}{All the fixed points of $f$ are in the region $\{z \in \mathbb{C} : \operatorname{Im}(z) > 0\}$}{There is no straight line in $\mathbb{C} \setminus \{3i\}$ which is mapped onto a straight line in $\mathbb{C}$ by $f$}
\ans{(C)}

\textbf{Q.35} The matrix $A = \begin{bmatrix} 1 & 2 & 0 \\ 1 & 3 & 1 \\ 0 & 1 & 3 \end{bmatrix}$ can be decomposed uniquely into the product $A = LU$, where $L = \begin{bmatrix} 1 & 0 & 0 \\ l_{21} & 1 & 0 \\ l_{31} & l_{32} & 1 \end{bmatrix}$ and $U = \begin{bmatrix} u_{11} & u_{12} & u_{13} \\ 0 & u_{22} & u_{23} \\ 0 & 0 & u_{33} \end{bmatrix}$. The solution of the system $LX = [1\;\;2\;\;2]^t$ is
\optsfourA{$[1\;\;1\;\;1]^t$}{$[1\;\;1\;\;0]^t$}{$[0\;\;1\;\;1]^t$}{$[1\;\;0\;\;1]^t$}
\ans{(A)}

\textbf{Q.36} Let $S = \left\{x \in \mathbb{R} : x \ge 0,\ \displaystyle\sum_{n=1}^{\infty} x^{\sqrt{n}} < \infty\right\}$. Then the supremum of $S$ is
\optsfourA{$1$}{$1/e$}{$0$}{$\infty$}
\ans{(A)}

\textbf{Q.37} The image of the region $\{z \in \mathbb{C} : \operatorname{Re}(z) > \operatorname{Im}(z) > 0\}$ under the mapping $z \mapsto e^{z^2}$ is
\optsfourA{$\{w \in \mathbb{C} : \operatorname{Re}(w) > 0,\ \operatorname{Im}(w) > 0\}$}{$\{w \in \mathbb{C} : \operatorname{Re}(w) > 0,\ \operatorname{Im}(w) > 0,\ |w| > 1\}$}{$\{w \in \mathbb{C} : |w| > 1\}$}{$\{w \in \mathbb{C} : \operatorname{Im}(w) > 0,\ |w| > 1\}$}
\ans{(C)}

\textbf{Q.38} Which of the following groups contains a unique normal subgroup of order four?
\optsfourA{$\mathbb{Z}_2 \oplus \mathbb{Z}_4$}{The dihedral group $D_4$ of order eight}{The quaternion group $Q_8$}{$\mathbb{Z}_2 \oplus \mathbb{Z}_2 \oplus \mathbb{Z}_2$}
\ans{(A) (Note: $\mathbb{Z}_2 \oplus \mathbb{Z}_4$ has 3 subgroups of order 4, $D_4$ has 3, $Q_8$ has 3, $\mathbb{Z}_2^3$ has 7; verify exact group in (A) against original paper, e.g. $\mathbb{Z}_8$ or $A_4$)}

\textbf{Q.39} Let $B$ be a real symmetric positive-definite $n \times n$ matrix. Consider the inner product on $\mathbb{R}^n$ defined by $\langle X, Y \rangle = Y^t B X$. Let $A$ be an $n \times n$ real matrix and let $T : \mathbb{R}^n \to \mathbb{R}^n$ be the linear operator defined by $T(X) = AX$ for all $X \in \mathbb{R}^n$. If $S$ is the adjoint of $T$, then $S(X) = CX$ for all $X \in \mathbb{R}^n$, where $C$ is the matrix
\optsfourA{$B^{-1} A^t B$}{$B A^t B^{-1}$}{$B^{-1} A B$}{$A^t$}
\ans{(A)}

\textbf{Q.40} Let $X$ be an arbitrary random variable that takes values in $\{0, 1, \dots, 10\}$. The minimum and maximum possible values of the variance of $X$ are
\optsfourA{$0$ and $30$}{$1$ and $30$}{$0$ and $25$}{$1$ and $25$}
\ans{(C)}

\textbf{Q.41} Let $M$ be the space of all $4 \times 3$ matrices with entries in the finite field of three elements. Then the number of matrices of rank three in $M$ is
\optsfourA{$(3^4 - 3)(3^4 - 3^2)(3^4 - 3^3)$}{$(3^4 - 1)(3^4 - 2)(3^4 - 3)$}{$(3^4 - 1)(3^4 - 3)(3^4 - 3^2)$}{$3^4(3^4 - 1)(3^4 - 2)$}
\ans{(C)}

\textbf{Q.42} Let $V$ be a vector space of dimension $m \ge 2$. Let $T : V \to V$ be a linear transformation such that $T^{n+1} = 0$ and $T^n \ne 0$ for some $n \ge 1$. Then which of the following is necessarily TRUE?
\optsfourA{$\operatorname{Rank}(T^n) \le \operatorname{Nullity}(T^n)$}{$\operatorname{trace}(T) \ne 0$}{$T$ is diagonalizable}{$n = m$}
\ans{(A)}

\textbf{Q.43} Let $X$ be a convex region in the plane bounded by straight lines. Let $X$ have $7$ vertices. Suppose $f(x, y) = ax + by + c$ has maximum value $M$ and minimum value $N$ on $X$ and $N < M$. Let $S = \{P : P \text{ is a vertex of } X \text{ and } N < f(P) < M\}$. If $S$ has $n$ elements, then which of the following statements is TRUE?
\optsfourA{$n$ cannot be $5$}{$n$ can be $2$}{$n$ cannot be $3$}{$n$ can be $4$}
\ans{(D)}

\textbf{Q.44} Which of the following statements are TRUE?\\
$P$: If $f \in L^1(\mathbb{R})$, then $f$ is continuous.\\
$Q$: If $f \in L^1(\mathbb{R})$ and $\lim_{|x|\to\infty} f(x)$ exists, then the limit is zero.\\
$R$: If $f \in L^1(\mathbb{R})$, then $f$ is bounded.\\
$S$: If $f \in L^1(\mathbb{R})$ is uniformly continuous, then $\lim_{|x|\to\infty} f(x)$ exists and equals zero.
\optsfourA{$Q$ and $S$ only}{$P$ and $R$ only}{$P$ and $Q$ only}{$R$ and $S$ only}
\ans{(A)}

\textbf{Q.45} Let $u$ be a real-valued harmonic function on $\mathbb{C}$. Let $g : \mathbb{R}^2 \to \mathbb{R}$ be defined by
\[
  g(x, y) = \int_0^{2\pi} u(e^{i\theta}(x + iy))\sin\theta\,d\theta.
\]
Which of the following statements is TRUE?
\optsfourA{$g$ is a harmonic polynomial}{$g$ is a polynomial but not harmonic}{$g$ is harmonic but not a polynomial}{$g$ is neither harmonic nor a polynomial}
\ans{(A)}

\textbf{Q.46} Let $S = \{z \in \mathbb{C} : |z| = 1\}$ with the induced topology from $\mathbb{C}$ and let $f : [0, 2] \to S$ be defined as $f(t) = e^{2\pi i t}$. Then which of the following is TRUE?
\optslistA{$K$ is closed in $[0, 2] \Rightarrow f(K)$ is closed in $S$}{$U$ is open in $[0, 2] \Rightarrow f(U)$ is open in $S$}{$f(X)$ is closed in $S \Rightarrow X$ is closed in $[0, 2]$}{$f(Y)$ is open in $S \Rightarrow Y$ is open in $[0, 2]$}
\ans{(A)}

\textbf{Q.47} Assume that all the zeros of the polynomial $a_n x^n + a_{n-1} x^{n-1} + \dots + a_1 x + a_0$ have negative real parts. If $u(t)$ is any solution to the ordinary differential equation
\[
  a_n \frac{d^n u}{dt^n} + a_{n-1}\frac{d^{n-1} u}{dt^{n-1}} + \dots + a_1 \frac{du}{dt} + a_0 u = 0,
\]
then $\displaystyle\lim_{t\to\infty} u(t)$ is equal to
\optsfourA{$0$}{$1$}{$n - 1$}{$\infty$}
\ans{(A)}

\section*{Common Data Questions}

\textbf{Common Data for Questions 48 and 49:}
Let $c_{00}$ be the vector space of all complex sequences having finitely many non-zero terms. Equip $c_{00}$ with the inner product $\langle x, y \rangle = \displaystyle\sum_{n=1}^{\infty} x_n \overline{y_n}$ for all $x = (x_n)$ and $y = (y_n)$ in $c_{00}$. Define $f : c_{00} \to \mathbb{C}$ by $f(x) = \displaystyle\sum_{n=1}^{\infty} \frac{x_n}{n}$. Let $N$ be the kernel of $f$.

\textbf{Q.48} Which of the following is FALSE?
\optslistA{$f$ is a continuous linear functional}{$\|f\| \le \dfrac{\pi}{\sqrt{6}}$}{There does not exist any $y \in c_{00}$ such that $f(x) = \langle x, y \rangle$ for all $x \in c_{00}$}{$N^\perp \ne \{0\}$}
\ans{(D)}

\textbf{Q.49} Which of the following is FALSE?
\optsfourA{$c_{00} \ne N$}{$N$ is closed}{$c_{00}$ is not a complete inner product space}{$c_{00} = N \oplus N^\perp$}
\ans{(D)}

\medskip
\textbf{Common Data for Questions 50 and 51:}
Let $X_1, X_2, \dots, X_n$ be an i.i.d.\ random sample from exponential distribution with mean $\mu$. In other words, they have density $f(x) = \begin{cases} \dfrac{1}{\mu}e^{-x/\mu}, & \text{if } x > 0, \\ 0, & \text{otherwise.} \end{cases}$

\textbf{Q.50} Which of the following is NOT an unbiased estimate of $\mu$?
\optsfourA{$X_1$}{$\dfrac{1}{n - 1}(X_2 + X_3 + \dots + X_n)$}{$n \cdot (\min\{X_1, X_2, \dots, X_n\})$}{$\dfrac{1}{n}\max\{X_1, X_2, \dots, X_n\}$}
\ans{(D)}

\textbf{Q.51} Consider the problem of estimating $\mu$. The m.s.e.\ (mean square error) of the estimate $T(X) = \dfrac{X_1 + X_2 + \dots + X_n}{n + 1}$ is
\optsfourA{$\mu^2$}{$\dfrac{1}{n + 1}\mu^2$}{$\dfrac{1}{(n + 1)^2}\mu^2$}{$\dfrac{n^2}{(n + 1)^2}\mu^2$}
\ans{(B)}

\section*{Linked Answer Questions}

\textbf{Statement for Linked Answer Questions 52 and 53:}
Let $X = \{(x, y) \in \mathbb{R}^2 : x^2 + y^2 = 1\} \cup ([-1, 1] \times \{0\}) \cup (\{0\} \times [-1, 1])$. Let $n_0 = \max\{k : k < \infty,\ \text{there are } k \text{ distinct points } p_1, \dots, p_k \in X \text{ such that } X \setminus \{p_1, \dots, p_k\} \text{ is connected}\}$.

\textbf{Q.52} The value of $n_0$ is \underline{\hspace{2cm}}.
\ans{$4$}

\textbf{Q.53} Let $q_1, \dots, q_{n_0 + 1}$ be $n_0 + 1$ distinct points and $Y = X \setminus \{q_1, \dots, q_{n_0 + 1}\}$. Let $m$ be the number of connected components of $Y$. The maximum possible value of $m$ is \underline{\hspace{2cm}}.
\ans{$8$}

\medskip
\textbf{Statement for Linked Answer Questions 54 and 55:}
Let $W(y_1, y_2)$ be the Wronskian of two linearly independent solutions $y_1$ and $y_2$ of the equation $y'' + P(x)y' + Q(x)y = 0$.

\textbf{Q.54} The product $W(y_1, y_2)P(x)$ equals
\optsfourA{$y_2 y_1'' - y_1 y_2''$}{$y_1 y_2' - y_2 y_1'$}{$y_1' y_2'' - y_2' y_1''$}{$y_2'' y_1' - y_1'' y_2'$}
\ans{(A)}

\textbf{Q.55} If $y_1 = e^{2x}$ and $y_2 = x e^{2x}$, then the value of $P(0)$ is
\optsfourA{$4$}{$-4$}{$2$}{$-2$}
\ans{(B)}

\section*{General Aptitude (GA) Questions}

\subsection*{Q.56 -- Q.60 Carry One Mark Each}

\textbf{Q.56} A number is as much greater than $75$ as it is smaller than $117$. The number is:
\optsfourA{$91$}{$93$}{$89$}{$96$}
\ans{(D)}

\textbf{Q.57} The professor \underline{ordered to} (I) \underline{the students} (II) \underline{to go} (III) \underline{out of the class} (IV).\\
Which of the above underlined parts of the sentence is grammatically incorrect?
\optsfourA{I}{II}{III}{IV}
\ans{(A)}

\textbf{Q.58} Which of the following options is the closest in meaning to the word given below:\\
\textbf{Primeval}
\optsfourA{Modern}{Historic}{Primitive}{Antique}
\ans{(C)}

\textbf{Q.59} Friendship, no matter how \underline{\hspace{2cm}} it is, has its limitations.
\optsfourA{cordial}{intimate}{secret}{pleasant}
\ans{(B)}

\textbf{Q.60} Select the pair that best expresses a relationship similar to that expressed in the pair:\\
\textbf{Medicine : Health}
\optsfourA{Science : Experiment}{Wealth : Peace}{Education : Knowledge}{Money : Happiness}
\ans{(C)}

\subsection*{Q.61 -- Q.65 Carry Two Marks Each}

\textbf{Q.61} $X$ and $Y$ are two positive real numbers such that $2X + Y \le 6$ and $X + 2Y \le 8$. For which of the following values of $(X, Y)$ will the function $f(X, Y) = 3X + 6Y$ give maximum value?
\optsfourA{$(4/3, 10/3)$}{$(8/3, 20/3)$}{$(8/3, 10/3)$}{$(4/3, 20/3)$}
\ans{(A)}

\textbf{Q.62} If $|4X - 7| = 5$, then the values of $2|X| - |-X|$ are:
\optsfourA{$2,\ 1/3$}{$1/2,\ 3$}{$3/2,\ 9$}{$2/3,\ 9$}
\ans{(B)}

\textbf{Q.63} The following table provides figures (in rupees) on annual expenditure of a firm for two years -- 2010 and 2011:
\begin{center}
\begin{tabular}{lcc}
  \toprule
  Category & 2010 & 2011 \\
  \midrule
  Raw material & 5200 & 6240 \\
  Power \& fuel & 7000 & 9450 \\
  Salary \& wages & 9000 & 12600 \\
  Plant \& machinery & 20000 & 25000 \\
  Advertising & 15000 & 19500 \\
  Research \& Development & 22000 & 26400 \\
  \bottomrule
\end{tabular}
\end{center}
In 2011, which of the following two categories have registered increase by same percentage?
\optslistA{Raw material and Salary \& wages}{Salary \& wages and Advertising}{Power \& fuel and Advertising}{Raw material and Research \& Development}
\ans{(D)}

\textbf{Q.64} A firm is selling its product at Rs.\ $60$ per unit. The total cost of production is Rs.\ $100$ and firm is earning total profit of Rs.\ $500$. Later, the total cost increased by $30\%$. By what percentage should the price be increased to maintain the same profit level?
\optsfourA{$5$}{$10$}{$15$}{$30$}
\ans{(A)}

\textbf{Q.65} Abhishek is elder to Savar. Savar is younger to Anshul.\\
Which of the given conclusions is logically valid and is inferred from the above statements?
\optslistA{Abhishek is elder to Anshul}{Anshul is elder to Abhishek}{Abhishek and Anshul are of the same age}{No conclusion follows}
\ans{(D)}

\newpage

% ============================================================================
% GATE 2014 -- GENERAL APTITUDE (GA) & MATHEMATICS (MA)
% ============================================================================

\begin{center}
  {\Large\bfseries GATE 2014 -- Mathematics (MA)}\\[4pt]
  {\small Duration: 180 minutes \quad\textbar\quad Maximum Marks: 100}
\end{center}

\section*{General Aptitude (GA) -- Set 1}

\subsection*{Q.1 -- Q.5 Carry One Mark Each}

\textbf{Q.1} A student is required to demonstrate a high level of \underline{comprehension} of the subject, especially in the social sciences.\\
The word closest in meaning to \textit{comprehension} is
\optsfourA{understanding}{meaning}{concentration}{stability}
\ans{(A)}

\textbf{Q.2} Choose the most appropriate word from the options given below to complete the following sentence.\\
One of his biggest \underline{\hspace{2cm}} was his ability to forgive.
\optsfourA{vice}{virtues}{choices}{strength}
\ans{(B)}

\textbf{Q.3} Rajan was not happy that Sajan decided to do the project on his own. On observing his unhappiness, Sajan explained to Rajan that he preferred to work independently.\\
Which one of the statements below is logically valid and can be inferred from the above sentences?
\optslistA{Rajan has decided to work only in a group.}{Rajan and Sajan were formed into a group against their wishes.}{Sajan had decided to give in to Rajan's request to work with him.}{Rajan had believed that Sajan and he would be working together.}
\ans{(D)}

\textbf{Q.4} If $y = 5x^2 + 3$, then the tangent at $x = 0$, $y = 3$
\optsfourA{passes through $x = 0, y = 0$}{has a slope of $+1$}{is parallel to the $x$-axis}{has a slope of $-1$}
\ans{(C)}

\textbf{Q.5} A foundry has a fixed daily cost of Rs.\ $50,000$ whenever it operates and a variable cost of Rs.\ $800Q$, where $Q$ is the daily production in tonnes. What is the cost of production in Rs.\ per tonne for a daily production of $100$ tonnes? \underline{\hspace{2cm}}
\ans{$1300$}

\subsection*{Q.6 -- Q.10 Carry Two Marks Each}

\textbf{Q.6} Find the odd one in the following group: \texttt{ALRVX}, \texttt{EPVZB}, \texttt{ITZDF}, \texttt{OYEIK}
\optsfourA{\texttt{ALRVX}}{\texttt{EPVZB}}{\texttt{ITZDF}}{\texttt{OYEIK}}
\ans{(D)}

\textbf{Q.7} Anuj, Bhola, Chandan, Dilip, Eswar, and Faisal live on different floors in a six-storeyed building (the ground floor is numbered $1$, the floor above it $2$, and so on). Anuj lives on an even-numbered floor. Bhola does not live on an odd-numbered floor. Chandan does not live on any of the floors below Faisal's floor. Dilip does not live on floor number $2$. Eswar does not live on a floor immediately above or immediately below Bhola. Faisal lives three floors above Dilip. Which of the following floor-person combinations is correct?
\begin{center}
\begin{tabular}{ccccccc}
  \toprule
  & Anuj & Bhola & Chandan & Dilip & Eswar & Faisal \\
  \midrule
  (A) & 6 & 2 & 5 & 1 & 3 & 4 \\
  (B) & 2 & 6 & 5 & 1 & 3 & 4 \\
  (C) & 4 & 2 & 6 & 3 & 1 & 5 \\
  (D) & 2 & 4 & 6 & 1 & 3 & 5 \\
  \bottomrule
\end{tabular}
\end{center}
\ans{(B)}

\textbf{Q.8} The smallest angle of a triangle is equal to two-thirds of the smallest angle of a quadrilateral. The ratio between the angles of the quadrilateral is $3 : 4 : 5 : 6$. The largest angle of the triangle is twice its smallest angle. What is the sum, in degrees, of the second largest angle of the triangle and the largest angle of the quadrilateral? \underline{\hspace{2cm}}
\ans{$180$}

\textbf{Q.9} One percent of the people of country $X$ are taller than $6$ ft. Two percent of the people of country $Y$ are taller than $6$ ft. There are thrice as many people in country $X$ as in country $Y$. Taking both countries together, what is the percentage of people taller than $6$ ft?
\optsfourA{$3.0$}{$2.5$}{$1.5$}{$1.25$}
\ans{(D)}

\textbf{Q.10} The monthly rainfall chart based on $50$ years of rainfall in Agra is shown in the original paper (displaying Average, $5$ percentile, and $95$ percentile rainfall in mm from Jan to Dec). Which of the following are true? ($k$ percentile is the value such that $k$ percent of the data fall below that value)\\
(i) On average, it rains more in July than in December\\
(ii) Every year, the amount of rainfall in August is more than that in January\\
(iii) July rainfall can be estimated with better confidence than February rainfall\\
(iv) In August, there is at least $500$ mm of rainfall
\optsfourA{(i) and (ii)}{(i) and (iii)}{(ii) and (iii)}{(iii) and (iv)}
\ans{(B)}

\section*{Mathematics (MA)}

\subsection*{Q.1 -- Q.25 Carry One Mark Each}

\textbf{Q.1} The function $f(z) = |z|^2 + i\bar{z} + 1$ is differentiable at
\optsfourA{$i$}{$1$}{$-i$}{no point in $\mathbb{C}$}
\ans{(C)}

\textbf{Q.2} The radius of convergence of the power series $\displaystyle\sum_{n=0}^{\infty} (-1)^n 2^{2n} z^{2n}$ is \underline{\hspace{2cm}}.
\ans{$0.5$ (official range: $0.49$ to $0.51$)}

\textbf{Q.3} Let $E_1$ and $E_2$ be two non-empty subsets of a normed linear space $X$ and let $E_1 + E_2 = \{x + y \in X : x \in E_1 \text{ and } y \in E_2\}$. Then which of the following statements is FALSE:
\optslistA{If $E_1$ and $E_2$ are convex, then $E_1 + E_2$ is convex}{If $E_1$ or $E_2$ is open, then $E_1 + E_2$ is open}{$E_1 + E_2$ must be closed if $E_1$ and $E_2$ are closed}{If $E_1$ is closed and $E_2$ is compact, then $E_1 + E_2$ is closed}
\ans{(C)}

\textbf{Q.4} Let $y(x)$ be the solution to the initial value problem $\dfrac{dy}{dx} = \sqrt{y + 2x}$ subject to $y(1.2) = 2$. Using the Euler method with the step size $h = 0.05$, the approximate value of $y(1.3)$, correct to two decimal places, is \underline{\hspace{2cm}}.
\ans{$2.21$ (official range: $2.2$ to $2.22$)}

\textbf{Q.5} Let $a \in \mathbb{R}$. If $ax$ is the polynomial which interpolates the function $f(x) = \sin x$ on $[-1, 1]$ at all the zeroes of the polynomial $4x^3 - 3x$, then $a$ is \underline{\hspace{2cm}}.
\ans{$0.47$ (official range: $0.46$ to $0.48$)}

\textbf{Q.6} If $u(x, t)$ is the D'Alembert's solution to the wave equation $\dfrac{\partial^2 u}{\partial t^2} = \dfrac{\partial^2 u}{\partial x^2}$, $x \in \mathbb{R}$, $t > 0$, with the condition $u(x, 0) = 0$ and $\dfrac{\partial u}{\partial t}(x, 0) = \cos x$, then $u\!\left(0, \dfrac{\pi}{4}\right)$ is \underline{\hspace{2cm}}.
\ans{$0.71$ (official range: $0.69$ to $0.72$)}

\textbf{Q.7} The solution to the integral equation $\varphi(x) = x + \displaystyle\int_0^x \sin(x - \xi)\varphi(\xi)\,d\xi$ is
\optsfourA{$x^2 + \dfrac{x^3}{3!}$}{$x - \dfrac{x^3}{3!}$}{$x + \dfrac{x^3}{3!}$}{$x^2 - \dfrac{x^3}{3!}$}
\ans{(C)}

\textbf{Q.8} The general solution to the ordinary differential equation $x^2\dfrac{d^2 y}{dx^2} + x\dfrac{dy}{dx} + \left(4x^2 - \dfrac{9}{25}\right)y = 0$ in terms of Bessel's functions $J_\nu(x)$ is
\optsfourA{$y(x) = c_1 J_{3/5}(2x) + c_2 J_{-3/5}(2x)$}{$y(x) = c_1 J_{3/10}(x) + c_2 J_{-3/10}(x)$}{$y(x) = c_1 J_{3/5}(x) + c_2 J_{-3/5}(x)$}{$y(x) = c_1 J_{3/10}(2x) + c_2 J_{-3/10}(2x)$}
\ans{(A)}

\textbf{Q.9} The inverse Laplace transform of $\dfrac{2s^2 - 4}{(s - 3)(s^2 - s - 2)}$ is
\optsfourA{$(1 + t)e^{-t} + \dfrac{7}{2}e^{-3t}$}{$\dfrac{e^t}{3} + te^{-t} + 2t$}{$\dfrac{7}{2}e^{3t} - \dfrac{e^{-t}}{6} - \dfrac{4}{3}e^{2t}$}{$\dfrac{7}{2}e^{-3t} - \dfrac{e^t}{6} - \dfrac{4}{3}e^{-2t}$}
\ans{(C)}

\textbf{Q.10} If $X_1, X_2$ is a random sample of size $2$ from an $N(0, 1)$ population, then $\dfrac{(X_1 + X_2)^2}{(X_1 - X_2)^2}$ follows
\optsfourA{$\chi_{(2)}^2$}{$F_{2, 2}$}{$F_{2, 1}$}{$F_{1, 1}$}
\ans{(D)}

\textbf{Q.11} Let $Z \sim N(0, 1)$ be a random variable. Then the value of $E[\max\{Z, 0\}]$ is
\optsfourA{$\dfrac{1}{\sqrt{\pi}}$}{$\sqrt{\dfrac{2}{\pi}}$}{$\dfrac{1}{\sqrt{2\pi}}$}{$\dfrac{1}{\pi}$}
\ans{(C)}

\textbf{Q.12} The number of non-isomorphic groups of order $10$ is \underline{\hspace{2cm}}.
\ans{$2$ (official range: $1.99$ to $2.01$)}

\textbf{Q.13} Let $a, b, c, d$ be real numbers with $a < c < d < b$. Consider the ring $C[a, b]$ with pointwise addition and multiplication. If $S = \{f \in C[a, b] : f(x) = 0 \text{ for all } x \in [c, d]\}$, then
\optslistA{$S$ is NOT an ideal of $C[a, b]$}{$S$ is an ideal of $C[a, b]$ but NOT a prime ideal of $C[a, b]$}{$S$ is a prime ideal of $C[a, b]$ but NOT a maximal ideal of $C[a, b]$}{$S$ is a maximal ideal of $C[a, b]$}
\ans{(B)}

\textbf{Q.14} Let $R$ be a ring. If $R[x]$ is a principal ideal domain, then $R$ is necessarily a
\optsfourA{Unique Factorization Domain}{Principal Ideal Domain}{Euclidean Domain}{Field}
\ans{(D)}

\textbf{Q.15} Consider the group homomorphism $\varphi : M_2(\mathbb{R}) \to \mathbb{R}$ given by $\varphi(A) = \operatorname{trace}(A)$. The kernel of $\varphi$ is isomorphic to which of the following groups?
\optsfourA{$M_2(\mathbb{R}) / \{A \in M_2(\mathbb{R}) : \varphi(A) = 0\}$}{$\mathbb{R}^2$}{$\mathbb{R}^3$}{$GL_2(\mathbb{R})$}
\ans{(C)}

\textbf{Q.16} Let $X$ be a set with at least two elements. Let $\tau$ and $\tau'$ be two topologies on $X$ such that $\tau' \ne \{\phi, X\}$. Which of the following conditions is necessary for the identity function $\mathrm{id} : (X, \tau) \to (X, \tau')$ to be continuous?
\optsfourA{$\tau \subseteq \tau'$}{$\tau' \subseteq \tau$}{no conditions on $\tau$ and $\tau'$}{$\tau \cap \tau' = \{\phi, X\}$}
\ans{(B)}

\textbf{Q.17} Let $A \in M_3(\mathbb{R})$ be such that $\det(A - I) = 0$, where $I$ denotes the $3 \times 3$ identity matrix. If $\operatorname{trace}(A) = 13$ and $\det(A) = 32$, then the sum of squares of the eigenvalues of $A$ is \underline{\hspace{2cm}}.
\ans{$81$ (official range: $80.99$ to $81.01$)}

\textbf{Q.18} Let $V$ denote the vector space $C^5[a, b]$ over $\mathbb{R}$ and $W = \left\{f \in V : \dfrac{d^4 f}{dt^4} + 2\dfrac{d^2 f}{dt^2} - f = 0\right\}$. Then
\optsfourA{$\dim(V) = \infty$ and $\dim(W) = \infty$}{$\dim(V) = \infty$ and $\dim(W) = 4$}{$\dim(V) = 6$ and $\dim(W) = 5$}{$\dim(V) = 5$ and $\dim(W) = 4$}
\ans{(B)}

\textbf{Q.19} Let $V$ be a real inner product space of dimension $10$. Let $x, y \in V$ be non-zero vectors such that $\langle x, y \rangle = 0$. Then the dimension of $\{x\}^\perp \cap \{y\}^\perp$ is \underline{\hspace{2cm}}.
\ans{$8$ (official range: $7.99$ to $8.01$)}

\textbf{Q.20} Consider the following linear programming problem: Minimize $x_1 + x_2$ subject to $2x_1 + x_2 \ge 8$, $2x_1 + 5x_2 \ge 10$, $x_1, x_2 \ge 0$. The optimal value to this problem is \underline{\hspace{2cm}}.
\ans{$4.25$ (official range: $4.24$ to $4.26$)}

\textbf{Q.21} Let $f(x) := \begin{cases} -3\pi, & \text{if } -\pi < x \le 0, \\ 3\pi, & \text{if } 0 < x < \pi, \end{cases}$ be a periodic function of period $2\pi$. The coefficient of $\sin 3x$ in the Fourier series expansion of $f(x)$ on the interval $[-\pi, \pi]$ is \underline{\hspace{2cm}}.
\ans{$4$ (official range: $3.99$ to $4.01$)}

\textbf{Q.22} For the sequence of functions $f_n(x) = \dfrac{1}{x^2}\sin\!\left(\dfrac{1}{nx}\right)$, $x \in [1, \infty)$, consider the following quantities expressed in terms of Lebesgue integrals:
\[
  I : \lim_{n\to\infty} \int_1^\infty f_n(x)\,dx, \qquad II : \int_1^\infty \lim_{n\to\infty} f_n(x)\,dx.
\]
Which of the following is TRUE?
\optslistA{The limit in $I$ does not exist}{The integrand in $II$ is not integrable on $[1, \infty)$}{Quantities $I$ and $II$ are well-defined, but $I \ne II$}{Quantities $I$ and $II$ are well-defined and $I = II$}
\ans{(D)}

\textbf{Q.23} Which of the following statements about the spaces $l^p$ and $L^p[0, 1]$ is TRUE?
\optsfourA{$l^3 \subset l^7$ and $L^6[0, 1] \subset L^9[0, 1]$}{$l^3 \subset l^7$ and $L^9[0, 1] \subset L^6[0, 1]$}{$l^7 \subset l^3$ and $L^6[0, 1] \subset L^9[0, 1]$}{$l^7 \subset l^3$ and $L^9[0, 1] \subset L^6[0, 1]$}
\ans{(B)}

\textbf{Q.24} The maximum modulus of $e^{z^2}$ on the set $S = \{z \in \mathbb{C} : 0 \le \operatorname{Re}(z) \le 1,\ 0 \le \operatorname{Im}(z) \le 1\}$ is
\optsfourA{$2/e$}{$e$}{$e + 1$}{$e^2$}
\ans{(B)}

\textbf{Q.25} Let $d_1, d_2$, and $d_3$ be metrics on a set $X$ with at least two elements. Which of the following is NOT a metric on $X$?
\optsfourA{$\min\{d_1, 2\}$}{$\max\{d_2, 2\}$}{$\dfrac{d_3}{1 + d_3}$}{$\dfrac{d_1 + d_2 + d_3}{3}$}
\ans{(B)}

\subsection*{Q.26 -- Q.55 Carry Two Marks Each}

\textbf{Q.26} Let $\Omega = \{z \in \mathbb{C} : \operatorname{Im}(z) > 0\}$ and let $C$ be a smooth curve lying in $\Omega$ with initial point $-1 + 2i$ and final point $1 + 2i$. The value of $\displaystyle\int_C \frac{1 + 2z}{1 + z}\,dz$ is
\optsfourA{$4 - \dfrac{1}{2}\ln 2 + i\dfrac{\pi}{4}$}{$-4 + \dfrac{1}{2}\ln 2 + i\dfrac{\pi}{4}$}{$4 + \dfrac{1}{2}\ln 2 - i\dfrac{\pi}{4}$}{$4 - \dfrac{1}{2}\ln 2 + i\dfrac{\pi}{2}$}
\ans{(A)}

\textbf{Q.27} If $a \in \mathbb{C}$ with $|a| < 1$, then the value of $\dfrac{1 - |a|^2}{\pi}\displaystyle\int_\Gamma \frac{|dz|}{|z + a|^2}$, where $\Gamma$ is the simple closed curve $|z| = 1$ taken with the positive orientation, is \underline{\hspace{2cm}}.
\ans{$2$ (official range: $1.99$ to $2.01$)}

\textbf{Q.28} Consider $C[-1, 1]$ equipped with the supremum norm given by $\|f\|_\infty = \sup\{|f(t)| : t \in [-1, 1]\}$ for $f \in C[-1, 1]$. Define a linear functional $T$ on $C[-1, 1]$ by $T(f) = \displaystyle\int_{-1}^{0} f(t)\,dt - \int_0^1 f(t)\,dt$ for all $f \in C[-1, 1]$. Then the value of $\|T\|$ is \underline{\hspace{2cm}}.
\ans{$2$ (official range: $1.99$ to $2.01$)}

\textbf{Q.29} Consider the vector space $C[0, 1]$ over $\mathbb{R}$. Consider the following statements:\\
$P$: If the set $\{t f_1, t^2 f_2, t^3 f_3\}$ is linearly independent, then the set $\{f_1, f_2, f_3\}$ is linearly independent, where $f_1, f_2, f_3 \in C[0, 1]$ and $t^n$ represents the polynomial function $t \mapsto t^n$, $n \in \mathbb{N}$.\\
$Q$: If $F : C[0, 1] \to \mathbb{R}$ is given by $F(x) = \displaystyle\int_0^1 x(t^2)\,dt$ for each $x \in C[0, 1]$, then $F$ is a linear map.\\
Which of the above statements hold TRUE?
\optsfourA{Only $P$}{Only $Q$}{Both $P$ and $Q$}{Neither $P$ nor $Q$}
\ans{(B)}

\textbf{Q.30} Using the Newton--Raphson method with the initial guess $x^{(0)} = 6$, the approximate value of the real root of $x\log_{10} x = 4.77$, after the second iteration, is \underline{\hspace{2cm}}.
\ans{$6.08$ (official range: $6.07$ to $6.09$)}

\textbf{Q.31} Let the following discrete data be obtained from a curve $y = y(x)$:
\begin{center}
\begin{tabular}{cccccc}
  \toprule
  $x$ & 0 & 0.25 & 0.5 & 0.75 & 1.0 \\
  \midrule
  $y$ & 1 & 0.9896 & 0.9589 & 0.9089 & 0.8415 \\
  \bottomrule
\end{tabular}
\end{center}
Let $S$ be the solid of revolution obtained by rotating the above curve about the $x$-axis between $x = 0$ and $x = 1$ and let $V$ denote its volume. The approximate value of $V$, obtained using Simpson's $\dfrac{1}{3}$ rule, is \underline{\hspace{2cm}}.
\ans{$2.82$ (official range: $2.81$ to $2.83$)}

\textbf{Q.32} The integral surface of the first order partial differential equation $2y(z - 3)\dfrac{\partial z}{\partial x} + (2x - z)\dfrac{\partial z}{\partial y} = y(2x - 3)$ passing through the curve $x^2 + y^2 = 2x$, $z = 0$ is
\optsfourA{$x^2 + y^2 - z^2 - 2x + 4z = 0$}{$x^2 + y^2 - z^2 - 2x + 8z = 0$}{$x^2 + y^2 + z^2 - 2x + 16z = 0$}{$x^2 + y^2 + z^2 - 2x + 8z = 0$}
\ans{(A)}

\textbf{Q.33} The boundary value problem $\dfrac{d^2\varphi}{dx^2} + \lambda\varphi = x$; $\varphi(0) = 0$ and $\dfrac{d\varphi}{dx}(1) = 0$, is converted into the integral equation $\varphi(x) = g(x) + \lambda \displaystyle\int_0^1 k(x, \xi)\varphi(\xi)\,d\xi$, where the kernel $k(x, \xi) = \begin{cases} \xi, & 0 < \xi < x, \\ x, & x < \xi < 1. \end{cases}$ Then $g\!\left(\dfrac{2}{3}\right)$ is \underline{\hspace{2cm}}.
\ans{$-0.28$ (official range: $-0.29$ to $-0.27$)}

\textbf{Q.34} If $y_1(x) = x$ is a solution to the differential equation $(1 - x^2)\dfrac{d^2 y}{dx^2} - 2x\dfrac{dy}{dx} + 2y = 0$, then its general solution is
\optsfourA{$y(x) = c_1 x + c_2(x\ln|1 + x^2| - 1)$}{$y(x) = c_1 x + c_2\!\left(\ln\left|\dfrac{1 - x}{1 + x}\right| + 1\right)$}{$y(x) = c_1 x + c_2\!\left(\dfrac{x}{2}\ln|1 - x^2| + 1\right)$}{$y(x) = c_1 x + c_2\!\left(\dfrac{x}{2}\ln\left|\dfrac{1 + x}{1 - x}\right| - 1\right)$}
\ans{(D)}

\textbf{Q.35} The solution to the initial value problem $\dfrac{d^2 y}{dt^2} + 2\dfrac{dy}{dt} + 5y = 3e^{-t}\sin t$, $y(0) = 0$ and $\dfrac{dy}{dt}(0) = 3$, is
\optsfourA{$y(t) = e^t(\sin t + \sin 2t)$}{$y(t) = e^{-t}(\sin t + \sin 2t)$}{$y(t) = 3e^t\sin t$}{$y(t) = 3e^{-t}\sin t$}
\ans{(B)}

\textbf{Q.36} The time to failure, in months, of light bulbs manufactured at two plants $A$ and $B$ obey the exponential distribution with means $6$ and $2$ months respectively. Plant $B$ produces four times as many bulbs as plant $A$ does. Bulbs from these plants are indistinguishable. They are mixed and sold together. Given that a bulb purchased at random is working after $12$ months, the probability that it was manufactured at plant $A$ is \underline{\hspace{2cm}}.
\ans{$0.93$ (official range: $0.92$ to $0.94$)}

\textbf{Q.37} Let $X, Y$ be continuous random variables with joint density function
\[
  f_{X,Y}(x, y) = \begin{cases} e^{-y}(1 - e^{-x}), & \text{if } 0 < x < y < \infty, \\ e^{-x}(1 - e^{-y}), & \text{if } 0 < y \le x < \infty. \end{cases}
\]
The value of $E[X + Y]$ is \underline{\hspace{2cm}}.
\ans{$4$ (official range: $3.99$ to $4.01$)}

\textbf{Q.38} Let $X = [0, 1) \cup (1, 2)$ be the subspace of $\mathbb{R}$, where $\mathbb{R}$ is equipped with the usual topology. Which of the following is FALSE?
\optslistA{There exists a non-constant continuous function $f : X \to \mathbb{Q}$}{$X$ is homeomorphic to $(-\infty, -3) \cup [0, \infty)$}{There exists an onto continuous function $f : [0, 1] \to \overline{X}$, where $\overline{X}$ is the closure of $X$ in $\mathbb{R}$}{There exists an onto continuous function $f : [0, 1] \to X$}
\ans{(D)}

\textbf{Q.39} Let $X = \begin{bmatrix} 2 & 0 & -3 \\ 3 & -1 & -3 \\ 0 & 0 & -1 \end{bmatrix}$. A matrix $P$ such that $P^{-1} X P$ is a diagonal matrix is
\optsfourA{$\begin{bmatrix} 1 & 1 & 1 \\ 0 & 1 & 1 \\ 1 & 1 & 0 \end{bmatrix}$}{$\begin{bmatrix} -1 & 1 & 1 \\ 0 & 1 & 1 \\ 1 & 1 & 0 \end{bmatrix}$}{$\begin{bmatrix} 1 & -1 & 1 \\ 0 & 1 & 1 \\ 1 & 1 & 0 \end{bmatrix}$}{$\begin{bmatrix} -1 & -1 & 1 \\ 0 & -1 & 1 \\ 1 & 1 & 0 \end{bmatrix}$}
\ans{(A)}

\textbf{Q.40} Using the Gauss--Seidel iteration method with the initial guess $\{x_1^{(0)} = 3.5,\ x_2^{(0)} = 2.25,\ x_3^{(0)} = 1.625\}$, the second approximation $\{x_1^{(2)}, x_2^{(2)}, x_3^{(2)}\}$ for the solution to the system of equations
\[
  2x_1 - x_2 = 7, \qquad -x_1 + 2x_2 - x_3 = 1, \qquad -x_2 + 2x_3 = 1,
\]
is
\optslistA{$x_1^{(2)} = 5.3125,\ x_2^{(2)} = 4.4491,\ x_3^{(2)} = 2.1563$}{$x_1^{(2)} = 5.3125,\ x_2^{(2)} = 4.3125,\ x_3^{(2)} = 2.6563$}{$x_1^{(2)} = 5.3125,\ x_2^{(2)} = 4.4491,\ x_3^{(2)} = 2.6563$}{$x_1^{(2)} = 5.4991,\ x_2^{(2)} = 4.4491,\ x_3^{(2)} = 2.1563$}
\ans{(B)}

\textbf{Q.41} The fourth order Runge--Kutta method given by $u_{j+1} = u_j + \dfrac{h}{6}[K_1 + 2K_2 + 2K_3 + K_4]$, $j = 0, 1, 2, \dots$, is used to solve the initial value problem $\dfrac{du}{dt} = u$, $u(0) = \alpha$. If $u(1) = 1$ is obtained by taking the step size $h = 1$, then the value of $K_4$ is \underline{\hspace{2cm}}.
\ans{$1.02$ (official range: $1.01$ to $1.03$)}

\textbf{Q.42} A particle $P$ of mass $m$ moves along the cycloid $x = (\theta - \sin\theta)$ and $y = (1 + \cos\theta)$, $0 \le \theta \le 2\pi$. Let $g$ denote the acceleration due to gravity. Neglecting the frictional force, the Lagrangian associated with the motion of the particle $P$ is:
\optsfourA{$m(1 - \cos\theta)\dot{\theta}^2 - mg(1 + \cos\theta)$}{$m(1 + \cos\theta)\dot{\theta}^2 + mg(1 + \cos\theta)$}{$m(1 + \cos\theta)\dot{\theta}^2 + mg(1 - \cos\theta)$}{$m(\theta - \sin\theta)\dot{\theta}^2 - mg(1 + \cos\theta)$}
\ans{(A)}

\textbf{Q.43} Suppose that $X$ is a population random variable with probability density function
\[
  f(x; \theta) = \begin{cases} \theta x^{\theta - 1}, & \text{if } 0 < x < 1, \\ 0, & \text{otherwise,} \end{cases}
\]
where $\theta$ is a parameter. In order to test the null hypothesis $H_0 : \theta = 2$ against the alternative hypothesis $H_1 : \theta = 3$, the following test is used: Reject the null hypothesis if $X_1 \ge \dfrac{1}{2}$ and accept otherwise, where $X_1$ is a random sample of size $1$ drawn from the above population. Then the power of the test is \underline{\hspace{2cm}}.
\ans{$0.875$ (official range: $0.86$ to $0.88$)}

\textbf{Q.44} Suppose that $X_1, X_2, \dots, X_n$ is a random sample of size $n$ drawn from a population with probability density function
\[
  f(x; \theta) = \begin{cases} \dfrac{x}{\theta^2}e^{-x/\theta}, & \text{if } x > 0, \\ 0, & \text{otherwise,} \end{cases}
\]
where $\theta$ is a parameter such that $\theta > 0$. The maximum likelihood estimator of $\theta$ is
\optsfourA{$\dfrac{\sum_{i=1}^n X_i}{n}$}{$\dfrac{\sum_{i=1}^n X_i}{n - 1}$}{$\dfrac{\sum_{i=1}^n X_i}{2n}$}{$\dfrac{2\sum_{i=1}^n X_i}{n}$}
\ans{(C)}

\textbf{Q.45} Let $\vec{F}$ be a vector field defined on $\mathbb{R}^2 \setminus \{(0, 0)\}$ by $\vec{F}(x, y) = \dfrac{y}{x^2 + y^2}\hat{i} - \dfrac{x}{x^2 + y^2}\hat{j}$. Let $\gamma, \alpha : [0, 1] \to \mathbb{R}^2$ be defined by $\gamma(t) = (8\cos 2\pi t,\, 17\sin 2\pi t)$ and $\alpha(t) = (26\cos 2\pi t,\, -10\sin 2\pi t)$. If $3\displaystyle\int_\alpha \vec{F} \cdot d\vec{r} - 4\int_\gamma \vec{F} \cdot d\vec{r} = 2m\pi$, then $m$ is \underline{\hspace{2cm}}.
\ans{$7$ (official range: $6.99$ to $7.01$)}

\textbf{Q.46} Let $g : \mathbb{R}^3 \to \mathbb{R}^3$ be defined by $g(x, y, z) = (3y + 4z,\, 2x - 3z,\, x + 3y)$ and let $S = \{(x, y, z) \in \mathbb{R}^3 : 0 \le x \le 1,\ 0 \le y \le 1,\ 0 \le z \le 1\}$. If $\displaystyle\iiint_{g(S)} (2x + y - 2z)\,dx\,dy\,dz = \alpha \iiint_S z\,dx\,dy\,dz$, then $\alpha$ is \underline{\hspace{2cm}}.
\ans{$75$ (official range: $74.99$ to $75.01$)}

\textbf{Q.47} Let $T_1, T_2 : \mathbb{R}^5 \to \mathbb{R}^3$ be linear transformations such that $\operatorname{rank}(T_1) = 3$ and $\operatorname{nullity}(T_2) = 3$. Let $T_3 : \mathbb{R}^3 \to \mathbb{R}^3$ be a linear transformation such that $T_3 \circ T_1 = T_2$. Then $\operatorname{rank}(T_3)$ is \underline{\hspace{2cm}}.
\ans{$2$ (official range: $1.99$ to $2.01$)}

\textbf{Q.48} Let $\mathbb{F}_3$ be the field of $3$ elements and let $\mathbb{F}_3 \times \mathbb{F}_3$ be the vector space over $\mathbb{F}_3$. The number of distinct linearly dependent sets of the form $\{u, v\}$, where $u, v \in (\mathbb{F}_3 \times \mathbb{F}_3) \setminus \{(0, 0)\}$ and $u \ne v$, is \underline{\hspace{2cm}}.
\ans{$4$ (official range: $3.99$ to $4.01$)}

\textbf{Q.49} Let $\mathbb{F}_{125}$ be the field of $125$ elements. The number of non-zero elements $\alpha \in \mathbb{F}_{125}$ such that $\alpha^5 = \alpha$ is \underline{\hspace{2cm}}.
\ans{$4$ (official range: $3.99$ to $4.01$)}

\textbf{Q.50} The value of $\displaystyle\iint_R xy\,dx\,dy$, where $R$ is the region in the first quadrant bounded by the curves $y = x^2$, $y + x = 2$, and $x = 0$, is \underline{\hspace{2cm}}.
\ans{$0.375$ (official range: $0.36$ to $0.38$)}

\textbf{Q.51} Consider the heat equation $\dfrac{\partial u}{\partial t} = \dfrac{\partial^2 u}{\partial x^2}$, $0 < x < \pi$, $t > 0$, with the boundary conditions $u(0, t) = 0$, $u(\pi, t) = 0$ for $t > 0$, and the initial condition $u(x, 0) = \sin x$. Then $u\!\left(\dfrac{\pi}{2}, 1\right)$ is \underline{\hspace{2cm}}.
\ans{$0.37$ (official range: $0.36$ to $0.38$)}

\textbf{Q.52} Consider the partial order in $\mathbb{R}^2$ given by the relation $(x_1, y_1) < (x_2, y_2)$ EITHER if $x_1 < x_2$ OR if $x_1 = x_2$ and $y_1 < y_2$. Then in the order topology on $\mathbb{R}^2$ defined by the above order,
\optslistA{$[0, 1] \times \{1\}$ is compact but $[0, 1] \times [0, 1]$ is NOT compact}{$[0, 1] \times [0, 1]$ is compact but $[0, 1] \times \{1\}$ is NOT compact}{both $[0, 1] \times [0, 1]$ and $[0, 1] \times \{1\}$ are compact}{both $[0, 1] \times [0, 1]$ and $[0, 1] \times \{1\}$ are NOT compact}
\ans{(D)}

\textbf{Q.53} Consider the following linear programming problem:
\begin{align*}
  \text{Minimize } & x_1 + x_2 + 2x_3 \\
  \text{Subject to } & x_1 + 2x_2 \ge 4, \\
                     & x_2 + 7x_3 \le 5, \\
                     & x_1 - 3x_2 + 5x_3 = 6, \\
                     & x_1, x_2 \ge 0,\quad x_3 \text{ is unrestricted.}
\end{align*}
The dual to this problem is:
\begin{align*}
  \text{Maximize } & 4y_1 + 5y_2 + 6y_3 \\
  \text{Subject to } & y_1 + y_3 \le 1, \\
                     & 2y_1 + y_2 - 3y_3 \le 1, \\
                     & 7y_2 + 5y_3 = 2,
\end{align*}
and further subject to:
\optsfourA{$y_1 \ge 0,\ y_2 \le 0$ and $y_3$ is unrestricted}{$y_1 \ge 0,\ y_2 \ge 0$ and $y_3$ is unrestricted}{$y_1 \ge 0,\ y_3 \le 0$ and $y_2$ is unrestricted}{$y_3 \ge 0,\ y_2 \le 0$ and $y_1$ is unrestricted}
\ans{(A)}

\textbf{Q.54} Let $X = C^1[0, 1]$. For each $f \in X$, define $p_1(f) := \sup\{|f(t)| : t \in [0, 1]\}$, $p_2(f) := \sup\{|f'(t)| : t \in [0, 1]\}$, and $p_3(f) := p_1(f) + p_2(f)$. Which of the following statements is TRUE?
\optsfourA{$(X, p_1)$ is a Banach space}{$(X, p_2)$ is a Banach space}{$(X, p_3)$ is NOT a Banach space}{$(X, p_3)$ does NOT have a denumerable basis}
\ans{(D)}

\textbf{Q.55} If the power series $\displaystyle\sum_{n=0}^{\infty} a_n (z + 3 - i)^n$ converges at $5i$ and diverges at $-3i$, then the power series
\optslistA{converges at $-2 + 5i$ and diverges at $2 - 3i$}{converges at $2 - 3i$ and diverges at $-2 + 5i$}{converges at both $2 - 3i$ and $-2 + 5i$}{diverges at both $2 - 3i$ and $-2 + 5i$}
\ans{(A)}

\end{document}
